% BEGIN INCLUDED SGA4_English_current_buildable_source_20260906/source/08/drafts/SGA4_VIII_LINES_43_576_English.texfrag
\setcounter{pageoriginal}{366}

% remarks:
% -- title too long at the top of the page

\section{平展拓扑斯的拓扑不变性}\etiquette[1]{VIII.sec1}

\begin{enonce*}{定理 1.1}\etiquette[1.1]{VIII.thm1.1}
设 \pageoriginale $f : X'\to X$ 是一个\emph{整的满射根态射}.
则基变换函子
$$
f^{\ast} : X_{\text{\rm ét}} \to X'_{\text{\rm ét}}
$$
是一个景的等价 (\ie 一个连续的范畴等价, 且其每个拟逆函子也连续).
因此, 函子
\begin{align*}
f_{\ast} : \widetilde{X'_{\text{\rm ét}}} &\to \widetilde{X_{\text{\rm ét}}}\\
f^{\ast} : \widetilde{X_{\text{\rm ét}}} &\to \widetilde{X'_{\text{\rm
      ét}}}
\end{align*}
是互为拟逆的范畴等价.
\end{enonce*}

当 $f$ 有限表示或 $f$ 平坦时, 第一个断言是众所周知的
(\SGA\sisi{}{ 1} IX 4.10 和 4.11);
一般情形归结到 $f$ 有限表示的情形. 事实上, 已知函子 $f^{\ast}$ 全忠实
(\SGA 1 IX 3.4). 只须证明 $f^{\ast}$\pageoriginale 本质满射,
\ie 每个 $X'$ 上的平展概形 $Z'$ 都“来自”某个 $X$ 上的平展概形 $Z$.
利用 $f$ 是普遍同胚以及 $f^{\ast}$ 全忠实这两个事实,
可归结到 $X$, $X'$ 和 $Z'$ 都是仿射概形的情形,
它们分别是环 $A$, $A'$ 和 $B'$ 的谱.
将 $A'$ 写成其有限生成 $A$-子代数 $A'_{i}$ 的归纳极限,
则 $B'$ 来自某个 $A'_{i}$ 上的平展代数 $B'_{i}$
(\cf \EGA IV 8), 这使我们归结到 $A'$ 在 $A$ 上整且\emph{有限型}的情形,
从而在 $A$ 上有限.
于是有一个 $A$-同构 $A'\simeq A''/J$,
其中 $A'' \simeq A [t_1, \ldots t_n]$, $J$ 是 $A''$ 的一个理想.
将 $J$ 写成其有限生成子理想 $J_i$ 的归纳极限, 并置 $A'_i = A'' /J_i$,
仍有 $A' = \varinjlim_i A'_{i}$,
所以 $B'$ 来自某个 $A'_i$ 上的平展代数 $B'_i$ (\emph{同上}).
另一方面, 由于 $\Spec A'\to \Spec A$ 是满射,
$\Spec A'_i \to \Spec A$ 也如此;
此外容易验证, 由于 $\Spec A'\to \Spec A$ 是根态射,
当 $i$ 充分大时 $\Spec A'_i \to \Spec A$ 也如此
(对 $\Spec A$ 应用 \EGA IV 1.9.9
\footnote{设 $S_i \subset S = \Spec A$ 是满足如下条件的 $s \in S$ 所成的集合:
$X_i = \Spec A'_i$ \sisi{作为集合}{在 $s$ 处}的纤维的可分秩为 1.
则 $S_i$ 是 $S$ 的\emph{可构造}子集 (\EGA IV 9.7.9),
它们构成一个递增族;
由关于 $S'=\Spec A'\to S$ 的假设,
以及 $S'_i$ 的纤维是 \emph{Noether} 概形这一事实, 它们的并为 $S$.
因此由 \EGA IV 1.9.9, 当 $i$ 充分大时 $S=S_i$, \cqfd.}).
这样就归结到 $A'$ 是某个 $A'_i$ 的情形, 从而它在 $A$ 上有限表示.

这就证明了
\sisi{\hyperref[VIII.prop2.1]{2.1}}{\hyperref[VIII.thm1.1]{1.1}}
的第一个断言. $f_{\ast}$ 和 $f^{\ast}$ 是互为拟逆的等价这一断言立即得出.

\begin{enonce*}{推论 1.2}\etiquette[1.2]{VIII.cor1.2}
设 $F$ 是 $X$ 上的一个 Abel 层, $F'$ 是它在 $X'$ 上的逆像.
则典范同态
\begin{equation*}
H^{i} (X, F) \to H^{i} (X', F')\tag{$\ast$}
\end{equation*}
都是同构. 同样, 若 $F'$ 是 $X'$ 上的一个 Abel 层,
$F$ 是它在 $X$ 上的正像, 则典范同态 $(\ast)$ 都是同构.
\end{enonce*}

\begin{enonce*}[remark]{例 1.3}\etiquette[1.3]{VIII.exem1.3}
下面\pageoriginale 是 \hyperref[VIII.thm1.1]{1.1} 的一些常用应用:
\begin{enumerate}
\renewcommand{\theenumi}{\alph{enumi}}
\renewcommand{\labelenumi}{\theenumi)}
\item $X'$ 是 $X$ 的一个子概形, 其底层集合与 $X$ 相同,
\ie 它由一个元素均幂零的理想 $I$ 定义.

\item $X$ 是可分闭域 $k$ 上的一个概形,
$k'$ 是 $k$ 的一个代数闭包, $X'=X\otimes_k k'$.

\item 设 $X$ 是一个几何单分支概形
(例如, 域 $k$ 上的一条代数曲线,
其奇点仅为“尖点”型奇点, 因而尤其排除了普通二重点).
则若 $X'$ 是 $X_{\red}$ 的正规化, 根据定义 $X'\to X$ 是根态射,
所以 \sisi{\hyperref[VIII.prop2.1]{2.1}}{\hyperref[VIII.thm1.1]{1.1}} 适用.
\end{enumerate}
\end{enonce*}

\begin{enonce*}[remark]{注 1.4}\etiquette[1.4]{VIII.rem1.4}
\hyperref[VIII.thm1.1]{1.1} 中的态射 $f$ 是一个\emph{普遍同胚},
\ie 它是同胚, 并且在每次基扩张后仍是同胚.
反过来, 若 $f$ 是一个普遍同胚,
可以证明 $f$ 满足
\sisi{\hyperref[VIII.prop2.1]{2.1}}{\hyperref[VIII.thm1.1]{1.1}}
的假设\footnote{若 $f$ 有限表示, \Cf \EGA IV 8.11.6;
一般情形见 \EGA IV 18.12.11.}, 这就解释了本节的标题.

关于例 \sisi{\hyperref[VIII.cor2.3]{2.3} c)}{\hyperref[VIII.exem1.3]{1.3} c)},
我们明确指出: 若 $X$ 是一条代数曲线
(为明确起见, 设在代数闭域 $k$ 上),
至少有一个不是“单分支”的奇点 (例如普通二重点),
则 (以 $X'$ 表示 $X_{\red}$ 的正规化)
\sisi{\hyperref[VIII.cor2.2]{2.2}}{\hyperref[VIII.cor1.2]{1.2}}
的结论即使对常系数的 $H^1$ 也不成立
(比较 \SGA 1 I 11 a) 和 \SGA 1 IX 5).
\end{enonce*}

\section{域的谱上的层}\etiquette[2]{VIII.sec2}\pageoriginale

\begin{enonce*}{命题 2.1}\etiquette[2.1]{VIII.prop2.1}
设 $k$ 是一个域, $\bar{k}$ 是 $k$ 的一个可分闭包,
$\pi = \Gal(\bar{k}/k)$ 是其拓扑 Galois 群,
$X=\Spec k$, $\bar{X}=\Spec\bar{k}$ ($\pi$ 在其上左作用).
考虑典范函子
$$
i : X_{\text{\rm ét}} \to \widetilde{X_{\text{\rm ét}}}
$$
(由 $i(X')(X'')=\Hom_X(X'',X')$ 定义), 以及典范函子
$$
j : \widetilde{X_{\text{\rm ét}}} \to \Bb_{\pi}
$$
它从 $X_{\text{\rm ét}}$ 上的层范畴到 $\pi$ 在其上连续左作用的
(离散) 集合组成的范畴 $\Bb_{\pi}$
\emph{(\cf IV \hyperref[IV.subsec2.7]{2.7})}, 由
$$
j(F) = \varinjlim_{\alpha} F(\Spec(k_{\alpha})) {\quoi},
$$
定义, 其中 $k_{\alpha}$ 遍历 $\bar{k}$ 的有限子扩张.
则 $i$ 和 $j$ 都是范畴等价.
(因此, 拓扑斯 $\widetilde{X_{\text{\rm ét}}}$
等价于 \emph{(IV \hyperref[IV.subsec3.4]{3.4})}
“分类拓扑斯” $\Bb_{\pi^{\bigdot}}$.)
\end{enonce*}

由 Galois 理论基本定理 (以 \SGA I V 的形式), 复合函子
$$
ji : X_{\text{\rm ét}} \to \widetilde{X_{\text{ét}}} \to (\pi\text{-}\Ens)
$$
是一个范畴等价. 由于 $\Bb_{\pi}$ 显然是一个拓扑斯 (II 4.14),
$X_{\text{ét}}$ 因此也如此.
此外, $X_{\text{ét}}$ 中的态射族 $(X_i\to X)$ 覆盖,
当且仅当它是满射的, \ie 它在 $\Bb_{\pi}$ 中的像是满射的,
或等价地, 是 $\Bb_{\pi}$ 的典范拓扑的覆盖族.
这表明 $X_{\text{ét}}$ 的拓扑确实是它的典范拓扑.
因此 $i$ 是一个等价, 从而 $j$ 也是.

\begin{enonce*}{推论 2.2}\etiquette[2.2]{VIII.cor2.2}
\hyperref[VIII.prop2.1]{2.1} 中的函子\pageoriginale $j$
诱导 $X=\Spec k$ 上的 Abel 层范畴\sisi{和}{与} Galois $\pi$-模范畴之间的一个等价.
\end{enonce*}

事实上, 后者恰是拓扑斯 $\Bb_{\pi}$ 中的“Abel 群”.
同样的论证, 通过“结构转移”, 给出:

\begin{enonce*}{推论 2.3}\etiquette[2.3]{VIII.cor2.3}
设 $F$ 是 $X=\Spec k$ 上的一个 Abel 层,
$M=j(F)$ 是相应的 Galois $\pi$-模.
则有一个关于 $F$ 的 $\partial$-函子的典范同构:
$$
H^{\ast} (X, F) \simeq H^{\ast} (\pi, M){\quoi},
$$
(其中右端表示 Galois 上同调, 例如在 Serre 的课程 \emph{C.G.} 中有研究).
\end{enonce*}

\begin{enonce*}{推论 2.4}\etiquette[2.4]{VIII.cor2.4}
假设 $k$ 可分闭. 则函子
$$
\Gamma : \widetilde{X_{\text{\rm ét}}} \to (\Ens)
$$
是一个范畴等价. 若 $F$ 是 $X=\Spec k$ 上的一个 Abel 层,
则当 $i\neq0$ 时 $H^{i}(X,F)=0$.
(换言之, 拓扑斯 $\widetilde{X_{\text{\rm ét}}}$ 是一个“点拓扑斯”
\emph{(IV \hyperref[IV.subsec2.2]{2.2})}.)
\end{enonce*}

\begin{enonce*}{推论 2.5}\etiquette[2.5]{VIII.cor2.5}
设 $k'$ 是可分闭域 $k$ 的一个可分闭扩张,
$X=\Spec k$, $X'=\Spec k'$, $u:X'\to X$ 是典范态射.
则函子
$$
u^{\ast} : \widetilde{X_{\text{\rm ét}}} \to \widetilde{X'_{\text{\rm ét}}}
$$
是一个范畴等价.
\end{enonce*}

\section{概形的几何点所对应的纤维函子}\etiquette[3]{VIII.sec3}\pageoriginale

\subsection{}\etiquette[3.1]{VIII.subsec3.1}

设 $X$ 是一个概形. 我们把每个作为\emph{可分闭}域 $\Omega$ 的谱的
$X$-概形 $\xi$ 称为 $X$ 的一个\emph{几何点}.
因此, 给定这样的点, 等价于给定一个普通点 $x\in X$,
再给定 $x$ 的剩余域 $k(x)$ 的一个\emph{可分闭}扩张 $\Omega$.
我们最常遇到的情形是 $\Omega$ 为 $k(x)$ 的一个\emph{可分闭包};
此时通常将它记为 $\bar{k(x)}$; 相应的 $X$ 的几何点记为 $\bar{x}$.
对给定的 $x\in X$, $\bar{k(x)}$ (分别 $\bar{x}$)
在相差\emph{不唯一}的 $k$-同构 (分别 $X$-同构) 的意义下确定.

\begin{enonce*}[remark]{注 3.2}\etiquette[3.2]{VIII.rem3.2}
在大多数几何问题中, 仅考虑 $\Omega$ \emph{代数闭}的情形更自然.
这里采用的不同约定是研究平展拓扑所特有的.
这正是我们在 \SGA 1 V 7 定义基本群时采用的约定.
不过注意, \emph{同上}的论述对于这里采用的约定也同样有效
(因为那里用到的 $\Omega$ 的性质是:
$\Omega$ 的谱的每个平展覆盖都是平凡的,
\ie 恰好是说 $\Omega$ 可分闭).
\end{enonce*}

\begin{enonce*}{定义 3.3}\etiquette[3.3]{VIII.def3.3}
设 $X$ 是一个概形, $\xi$ 是 $X$ 的一个几何点,
$u:\xi\to X$ 是它的结构态射.
对应于 $\xi$ 的 (几何) \emph{纤维函子}, 记为 $F\mapsto F_{\xi}$,
是如下的复合函子
$$
\widetilde{X_{\text{\rm ét}}} \to (\Ens)
$$
其分解为
$$
\widetilde{X_{\text{\rm ét}}} \xrightarrow{u^{\ast}}
\widetilde{\xi_{\text{\rm ét}}} \xrightarrow{\Gamma_{\xi}} (\Ens){\quoi}.
$$
\end{enonce*}

等价地, 根据 \hyperref[VIII.cor2.4]{2.4},
可以说 $X$ 的一个几何点 $\xi$ 定义了拓扑斯 $X^{\sim}_{\text{ét}}$
的一个点 (IV \hyperref[IV.def6.1]{6.1}),
我们取的就是与它相应的纤维函子 (\loccit).

\setcounter{subsection}{3}
\subsection{}\etiquette[3.4]{VIII.subsec3.4}

由于\pageoriginale 函子 $\Gamma_{\xi}$ ($\xi$ 上的整体截面函子)
是一个范畴\emph{等价} \hyperref[VIII.cor2.4]{(2.4)},
知道纤维函子 $F\mapsto F_{\xi}$,
实质上等价于知道逆像函子 $u^{\ast}$.
此外, 从 \hyperref[VIII.cor2.5]{2.5} 立即可知:
若 $\xi$ 是 $X$ 的一个几何点, 且存在一个 $X$-态射 $\xi'\to\xi$
(\ie $\xi'$ 对应于 $\Omega$ 的一个可分闭扩张 $\Omega'$),
则典范同态 $F_{\xi}\to F_{\xi'}$ 是一个\emph{函子同构}
$$
F_{\xi} \simeq F_{\xi'}{\quoi}.
$$
这表明, 为了研究纤维函子,
可以 (将 $\Omega$ 换成 $k$ 在 $\Omega$ 中的可分代数闭包后)
仅考虑 $\Omega = k(\xi)$ 是 $k(x)$ 的一个可分闭包 $\bar{k(x)}$ 的情形;
并且, 位于同一个 $x \in X$ 上方的 $X$ 的几何点所对应的纤维函子彼此同构
(非典范地).

按照 \hyperref[VIII.subsec3.1]{3.1} 的约定,
我们用 $F_{\bar{x}}$ 表示选定 $k(x)$ 的一个可分闭包
$\bar{k(x)}$ 后所对应的纤维函子.

我们指出一个显然的传递性性质
(它说明, 不把自己局限于剩余域的可分闭包所定义的几何点, 在技术上很有用):
设 $f : X'\to X$ 是一个概形态射,
$u' : \xi'\to X'$ 是 $X'$ 的一个几何点.
则 $u = f u': \xi'\to X$ 是 $X$ 的一个几何点, 将它记为 $\xi$;
对于这些几何点, 存在一个关于 $F \in \Ob \widetilde{X_{\text{ét}}}$
具有函子性的同构:
$$
f^{\ast} (F)_{\xi'} \simeq F_{\xi}{\quoi}.
$$
事实上, 这\pageoriginale 来自逆像函子的传递公式
$$
(fu')^{\ast} \xrightarrow{\sim} u'^{\ast} f^{\ast}{\quoi}.
$$

\begin{enonce*}{定理 3.5}\etiquette[3.5]{VIII.thm3.5}
设 $X$ 是一个概形.
\begin{enumerate}
\renewcommand{\theenumi}{\alph{enumi}}
\renewcommand{\labelenumi}{\theenumi)}
\item 对 $X$ 的每个几何点 $\xi$,
纤维函子 $F \mapsto F_{\xi}$,
$\widetilde{X_{\text{\rm ét}}} \to (\Ens)$
与任意 $\varinjlim$ (以范畴 $\in \Uu$ 为指标)
以及有限 $\varprojlim$ 交换\footnote{因此它确实是第 IV 讲意义下的一个“纤维函子”.}.
特别地, 它将群层 (\resp Abel 层, \etc)
映为群 (\resp Abel 群, \etc).

\item 当 $x$ 遍历 $X$ 的点时,
纤维函子 $F_{\bar{x}}$ 构成一个“保守”函子族,
\ie 若 $u : F \to G$ 是 $X$ 上的一个层同态,
则 $u$ 是同构, 当且仅当对每个 $x \in X$,
相应同态 $u_{\bar{x}} : F_{\bar{x}} \to G_{\bar{x}}$ 都是同构
\footnote{因此, 用 IV \hyperref[IV.prop6.5]{6.5} 的术语来说,
形如 $F \mapsto F_{\bar{x}}$ 的纤维函子“足够多”.}.
\end{enumerate}
\end{enonce*}

考虑到 a) 中的正合性性质, 我们立即指出 b) 的以下形式推论:

\begin{enonce*}{推论 3.6}\etiquette[3.6]{VIII.cor3.6}
设 $u : F \to G$ 是一个层同态.
则 $u$ 是单态射 (\resp 满态射), 当且仅当对每个 $x \in X$,
$u_{\bar{x}}:F_{\bar{x}}\to G_{\bar{x}}$ 是单射 (\resp 满射).
\end{enonce*}

\begin{enonce*}{推论 3.7}\etiquette[3.7]{VIII.cor3.7}
设 $u$, $v : F \to G$ 是 $X$ 上的两个层态射.
则 $u = v$, 当且仅当对每个 $x \in X$,
都有 $u_{\bar{x}} = v_{\bar{x}} : F_{\bar{x}} \to G_{\bar{x}}$.
特别地, 若 $u$, $v$ 是 $F$ 的两个截面, 则\pageoriginale
$u = v$, 当且仅当对每个 $x \in X$,
在 $F_{\bar{x}}$ 中都有 $u_{\bar{x}} = v_{\bar{x}}$.
\end{enonce*}

\begin{enonce*}{推论 3.8}\etiquette[3.8]{VIII.cor3.8}
设 $F \to G \to H$ 是 $X$ 上两个层同态组成的序列.
则这个序列正合, 当且仅当对每个 $x \in X$,
相应序列 $F_{\bar{x}} \to G_{\bar{x}} \to H_{\bar{x}}$ 正合.
\end{enonce*}

这些推论的推导留给读者\footnote{\Cf IV \hyperref[IV.exer6.5.2]{6.5.2}.},
%% issue: check that (**) is indeed present upon compilation
不过, 以下证明也会证明其中的一部分.

断言 a) 来自 (\hyperrefext[VII][VII.subsec1.4]{1.4}) 中已经指出的正合性性质,
它适用于每个逆像函子, 例如
$u^{\ast} : \widetilde{X_{\text{ét}}} \to
\widetilde{\xi_{\text{ét}}}$,
再结合 $\widetilde{\xi_{\text{ét}}} \to (\Ens)$ 是一个等价这一事实即可.
为证明 b), 我们指出以下公式, 它给出了\emph{计算}纤维函子的方法:

\begin{enonce*}{命题 3.9\footnotemark}\etiquette[3.9]{VIII.prop3.9}
\footnotetext{\Cf IV \hyperref[IV.eq6.8.3]{6.8.3}} %% issue: check the number of * upon compilation
设 $\xi \xrightarrow{u} X$ 是 $X$ 的一个几何点,
$C_{\xi}$ 是“在 $\xi$ 处带点的平展 $X$-概形”范畴,
\ie 其对象为二元组 $(X', u')$,
其中 $X'\in \Ob X_{\text{\rm ét}}$ 是 $X$ 上的平展概形,
$u'$ 是一个 $X$-态射 $\xi \to X'$.
则反范畴 $C^{0}_{\xi}$ 滤过, 并且当 $F$
(\resp 预层 $P$) 在 $X$ 上的层中变化时, 有一个典范函子同构
$$
F_{\xi} \simeq \varinjlim_{C^{0}_{\xi}} F(X'), (\emph{\resp}
(aP)_{\xi} \simeq \varinjlim_{C^{0}_{\xi}} P(X')){\quoi},
$$
(其中 $aP$ = 与 $P$ 相伴的层).
\end{enonce*}

事实上, 先注意: 对 $\xi_{\text{ét}}$ 上的每个预层 $P$,
若 $aP$ 是相伴层, 则自然同态
$$
P(\xi) \to aP(\xi)
$$
是一个同构, 例如这可直接由将\pageoriginale $aP$ 显式计算为 $L(LP)$
(II \hyperref[II.theo3.4]{3.19}) 得出.
设 $\varphi : \widetilde{X_{\text{ét}}} \to \widetilde{\xi_{\text{ét}}}$
为“逆像”函子. 则 $u^{\ast} = \varphi^{\ast}$
正是 (III \hyperref[III.prop2.3]{2.3}. (d)) 复合 $a\varphi^{\bigdot} i$,
其中 $i : \widetilde{X_{\text{ét}}} \to \widehat{X_{\text{ét}}}$
是包含函子, $a : \widehat{\xi_{\text{ét}}} \to
\widetilde{\xi_{\text{ét}}}$ 是“相伴层”函子.
由前面的观察, 得到
$$
u^{\ast} (F) (\xi) = (\varphi^{\bigdot} i(F)) (\xi){\quoi},
$$
所以只须应用有关证明中给出的 $\varphi^{\bigdot}$ 的显式表达式.
\SGAArchiveOnly{\nde{冻结源稿没有指明准确的内部引用.}}
若从 $X$ 上的一个预层 $P$ 出发, 相应的情形以同样方式证明,
只须利用函子同构 $\varphi^{\ast} a(P) \simeq a
\varphi^{\ast} (P)$.
\SGAArchiveOnly{\nde{冻结源稿没有指明准确的内部引用.}}
范畴 $C_{\xi}$ 滤过这一事实
(它对\emph{每个} $X$-概形 $\xi$ 都成立, 不要求其退化为一个几何点),
直接来自 $C_{\xi}$ 中存在有限投射极限这一事实,
后者又来自 $X_{\text{ét}}$ 是 $(\Sch) /X$ 的一个对有限投射极限稳定的子范畴.

现在设 $u : F \to G$ 是 $X$ 上的一个层同态,
满足对每个 $x \in X$, $u_{\bar{x}} : F_{\bar{x}} \to G_{\bar{x}}$ 是单态射;
我们证明 $u$ 是单态射.
为此须证明: 若 $X'\in \Ob X_{\text{ét}}$,
$\varphi$, $\psi \in F(X')$ 满足 $u (\varphi) = u (\psi)$,
则 $\varphi = \psi$.
将 $X$ 换成 $X'$ 并利用纤维函子的传递性,
可假设 $X = X'$.
于是对每个 $x \in X$, 有 $u (\varphi)_{\bar{x}} = u (\psi)_{\bar{x}}$,
\ie $u_{\bar{x}} (\varphi_{\bar{x}}) = u_{\bar{x}} (\psi_{\bar{x}})$,
从而 $\varphi_{\bar{x}} = \psi_{\bar{x}}$,
因为 $u_{\bar{x}}$ 是单态射.
利用 \hyperref[VIII.prop3.9]{3.9}, %% XXX Reference to correct}
可知存在一个 $X'_x \in \Ob X_{\text{ét}}$, 其像包含 $x$,
使 $\varphi$ 和 $\psi$ 在 $X'_x$ 上的逆像相同.
由于当 $x$ 在 $X$ 中变化时, $X'_x$ 构成 $X$ 的一个覆盖族,
可得 $\varphi = \psi$.

这就\pageoriginale 证明了: 若 $u_{\bar{x}}$ 都是单态射, 则 $u$ 也是单态射.
再假设 $u_{\bar{x}}$ 都是满态射; 我们证明 $u$ 也如此, 从而 $u$ 是同构.
只须证明: 对每个 $X'\in \Ob X_{\text{ét}}$ 和每个 $\psi \in G (X')$,
存在 $\varphi \in F(X')$ 使 $u(X') (\varphi) = \psi$.
仍可假设 $X'= X$, 再利用 \hyperref[VIII.prop3.9]{3.9}, %% XXX Reference to correct
可见对每个 $x \in X$, \sisi{}{都}存在 $X'_x \in \Ob X_{\text{ét}}$,
使 $X'_x$ 在 $X$ 中的像包含 $x$,
并存在 $\varphi_x \in F(X'_x)$,
其在 $G (X'_x)$ 中的像是 $\psi$ 的逆像.
利用 $u$ 是单态射这一事实,
可见 $\varphi_{x}$ 在 $\fprod{X'_x}{X'_y}{X}$ $(x, y \in X)$ 上相容,
因而来自唯一确定的一个 $\varphi \in F (X)$;
此时 $u(\varphi) = \psi$, 因为它们在 $X'_{x}$ 上的逆像相同. \cqfd.

\begin{enonce*}[remark]{附注 3.10}\etiquette[3.10]{VIII.sch3.10}
定理 \hyperref[VIII.thm3.5]{3.5} 表明,
平展拓扑的“纤维函子”全体具有与拓扑空间上的通常层论相同的主要性质.
我们将非常频繁地使用 \hyperref[VIII.thm3.5]{3.5} 及其推论,
因此\sisi{}{将}一般不再明确援引它们.
\end{enonce*}

\setcounter{subsection}{10}
\subsection{}\etiquette[3.11]{VIII.subsec3.11}

对每个 $x \in X$, 我们也将滥用语言,
用 $x$ 表示带有通常结构态射
$$
i_x : x = \Spec k(x) \to X{\quoi}.
$$
的 $X$-概形 $\Spec k(x)$.
它给出一个典范函子
$$
i_x^{\ast} : \widetilde{X_{\text{ét}}} \to \widetilde{x_{\text{ét}}}{\quoi}.
$$
若 $F$ 是 $X$ 上的一个层, 则置 $F_x = i^{\ast}_x (F)$;
因此它是 $x = \Spec k(x)$ 上的一个层
(从而根据 \hyperref[VIII.prop2.1]{2.1},
也可将它与 $x$ 上的一个\pageoriginale 平展概形等同).
它关于 $F$ 具有函子性, 而由 \hyperrefext[VII][VII.subsec1.4]{1.4},
函子 $i^{\ast}_x : F \mapsto F_x$ 仍与任意 $\varinjlim$ 和有限 $\varprojlim$ 交换.

若 $\bar{x} = \Spec \bar{k(x)}$,
则纤维函子 $F \mapsto F_{\bar{x}}$ 典范同构于函子 $F \mapsto F_x$
与 \sisi{\hyperref[VIII.thm1.1]{1.1}}{\hyperref[VIII.prop2.1]{2.1}}
中的函子 $j$ 的复合
(这里用到后者同构于 $\Spec k$ 的几何点 $\Spec \bar{k}$ 所对应的纤维函子这一事实).
由此及 \hyperref[VIII.thm3.5]{3.5} 立即可知,
函子系 $(F \mapsto F_x)$ $(x \in X)$ 同样保守.
此外, 若 $k(x)$ 可分闭, \ie $x = \bar{x}$,
则严格说来记号 $F_x$ 和 $F_{\bar{x}}$ 互相冲突
(前者表示 $\widetilde{x_{\text{ét}}}$ 的一个对象, 后者表示一个集合),
不过根据 \sisi{\hyperref[VIII.thm1.1]{1.1}}{\hyperref[VIII.prop2.1]{2.1}},
这种冲突无关紧要.

\subsection{}\etiquette[3.12]{VIII.subsec3.12}

前面的观察表明, 若 $\bar{x} = \Spec \bar{k(x)}$,
则对 $X$ 上的每个层 $F$, 群 $\pi_x = \Gal (\bar{k(x)}/ k(x))$
自然地 (从左边) 作用在纤维集合 $F_{\bar{x}}$ 上,
所以实际上可将 $F \mapsto F_{\bar{x}}$ 看成一个函子
$$
\widetilde{X_{\text{ét}}} \to \Bb_{\pi_x}
$$
根据 \sisi{\hyperref[VIII.thm1.1]{1.1}}{\hyperref[VIII.prop2.1]{2.1}},
知道这个函子实质上等价于知道函子 $F \mapsto F_x$.

这在一个重要的特例中表明,
与拓扑空间上的层的情形不同,“纤维函子”一般具有非平凡的自同构.
在一般情形中, 我们将在 \hyperref[VIII.thm7.9]{7.9}
确定一个纤维函子到另一个纤维函子的全部同态.

\begin{enonce*}[remark]{注 3.13}\etiquette[3.13]{VIII.rem3.13}
下面\pageoriginale 是 \hyperref[VIII.thm3.5]{3.5} b) 的一个有时有用的推广.
考虑 $X$ 的一个子集 $E$, 满足对每个平展态射 $f : X' \to X$,
$E' = f^{-1} (E)$ 在 $X'$ 中“很稠密” (EGA IV 10.1.3),
\ie 对上述每个 $f$, $X'$ 的每个包含 $E'$ 的开子集 $U$ 都等于 $X'$.
则当 $x \in E$ 时, $\widetilde{X_{\text{\rm ét}}}$ 上的纤维函子 $F_{\bar{x}}$
已经构成一个“保守族”
(因此, 在推论 \hyperref[VIII.cor3.6]{3.6}, \hyperref[VIII.cor3.7]{3.7}
和 \hyperref[VIII.cor3.8]{3.8} 中, 只选取 $x \in E$ 时, 它们仍成立).
重复 \hyperref[VIII.thm3.5]{3.5} b) 的证明, 立即得出这一点.
上述结果尤其适用于以下情形:
\begin{enumerate}
\renewcommand{\theenumi}{\alph{enumi}}
\renewcommand{\labelenumi}{\theenumi)}
\item $X$ 是一个 Jacobson 概形 (EGA IV 10.4.1.),
例如在某个域或 $\Spec(𝒵)$ 上局部有限型,
$E$ 是 $X$ 的闭点集.

\item $X$ 在某个概形 $S$ 上局部有限型,
$E$ 是 $X$ 中在其纤维内为闭点的点所成的集合.

\item $X$ 局部 Noether,
$E$ 是满足 $\bar{x}$ 为有限集 (\ie 维数 $\leq 1$ 的半局部概形)
的 $x \in X$ 所成的集合,
\cf \EGA IV 10.5.3 和 10.5.5.
\end{enumerate}
\end{enonce*}

% END INCLUDED SGA4_English_current_buildable_source_20260906/source/08/drafts/SGA4_VIII_LINES_43_576_English.texfrag


% BEGIN INCLUDED SGA4_English_current_buildable_source_20260906/source/08/drafts/SGA4_VIII_LINES_577_804_English.texfrag
\section{严格局部环与严格局部概形}\etiquette[4]{VIII.sec4}

\subsection{}\etiquette[4.1]{VIII.subsec4.1}

回顾一下, 按照 Azumaya,
若局部环 $A$ 满足以下等价条件, 就称它是 \emph{Hensel 环}:
\begin{enumerate}
\renewcommand{\theenumi}{\roman{enumi}}
\renewcommand{\labelenumi}{(\theenumi)}
\item 每个有限 $A$-代数 $B$ 都是局部环的直积 $\prod_i B_i$
(此时 $B_i$ 必同构于 $B_{\mathfrak m_i}$,
其中\pageoriginale $\mathfrak m_i$ 遍历 $B$ 的极大理想).

\item 若一个 $A$-代数 $B$ 是有限型 $A$-代数 $C$ 在某个素理想 $\mathfrak p$
处的局部化, 其中该素理想位于 $A$ 的极大理想 $\mathfrak m$ 之上,
并且 $B/\mathfrak mB$ 在 $k=A/\mathfrak m$ 上有限,
则它在 $A$ 上有限.

\item 与 (ii) 相同, 但假设 $C$ 在 $\mathfrak p$ 处在 $A$ 上平展,
且 $k\to B/\mathfrak mB$ 是同构.

\item 经典形式的“ Hensel 引理”对 $A$ 成立,
即对每个首一多项式 $F\in A[T]$,
以 $\overset{\bigdot}{F}\in k[T]$ 表示相应的约化多项式,
则对 $\overset{\bigdot}{F}$ 的每个分解
$$
\overset{\bigdot}{F}
=\overset{\bigdot}{F}_1\overset{\bigdot}{F}_2
$$
其中两个因子是互素首一多项式,
$\overset{\bigdot}{F}_1$ 和 $\overset{\bigdot}{F}_2$
都可提升为 $F_1,F_2\in A[T]$, 使 $F=F_1F_2$ (提升必唯一).

关于这些条件的等价性证明以及 Hensel 环的一般性质, 见 \EGA IV 18.5.11.
回顾一下, 对任意局部环 $A$, 可以证明 (Nagata), 存在一个局部同态
$$
A\to A^h
$$
其中 $A^h$ 是 \emph{Hensel 环},
而且这个同态对于从 $A$ 到局部 \sisi{Noether}{Hensel} 环的局部同态是\emph{泛的}.
环 $A^h$ 称为 $A$ 的 Hensel 化.
\EGA IV 18.6 中给出的构造 (与 Nagata 的不同) 是令
$$
A^h=\varinjlim_i A_i{\quoi},
$$
其中 $A_i$ 遍历在 $A$ 上本质有限型且本质平展\pageoriginale
(即某个平展 $A$-代数的局部化) 的 $A$-代数,
满足 $A\to A_i$ 是局部同态, 且剩余扩张 $k(A_i)/k(A)$ 平凡.
这个构造尤其表明 $A^h$ 在 $A$ 上平坦,
并且若 $A$ 是 Noether 环, 它的 Hensel 化也是.
\end{enumerate}

\begin{enonce*}{定义 4.2}\etiquette[4.2]{VIII.def4.2}
若局部环 $A$ \emph{满足以下等价条件之一}, 就称它严格局部:
\begin{enumerate}
\renewcommand{\theenumi}{\roman{enumi}}
\renewcommand{\labelenumi}{(\theenumi)}
\item $A$ 是 Hensel 环 (\hyperref[VIII.subsec4.1]{4.1}),
且其剩余域可分闭.

\item 每个局部同态 $A\to B$ 都是同构,
其中 $B$ 是 $A$ 上某个平展代数的局部化.
\end{enumerate}

若概形同构于某个严格局部环的谱, 就称它是\emph{严格局部概形}.
\end{enonce*}

(i) 和 (ii) 的等价性直接由 \hyperref[VIII.subsec4.1]{4.1} 得出.
注意, 用几何语言表述,
\hyperref[VIII.def4.2]{4.2} (ii)
(分别 \hyperref[VIII.subsec4.1]{4.1} (iii)) 是说:
对每个平展 $Y$-概形 $X$ 以及纤维 $X_y$ 的每个点 $x$
(分别 $X_y$ 的每个在 $k(y)$ 上\emph{有理}的点),
其中 $y$ 是 $Y$ 的闭点,
存在 $X$ 在 $Y$ 上的一个经过 $x$ 的截面 (且此截面必唯一).

\begin{enonce*}{定义 4.3}\etiquette[4.3]{VIII.def4.3}
设 $X$ 是一个概形, $\Oo_{X_{\text{\rm ét}}}$ 是 $X_{\text{\rm ét}}$ 上由
$$
\Oo_{X_{\text{\rm ét}}}(X')=\Gamma(X',\Oo_{X'})
$$
定义的层 (比较 \emph{\hyperrefext[VII][VII.sec2]{2}. c)}),
$u:\xi=\Spec(\Omega)\to X$ 是 $X$ 的一个几何点.
\emph{$X$ 在 $\xi$ 处的严格局部环}, 记为
$\Oo_{{X,}_{\xi}}$ 或简记\pageoriginale 为 $\Oo_\xi$,
是层 $\Oo_{X_{\text{\rm ét}}}$ 在点 $\xi$ 处的纤维.
它的谱称为\emph{$X$ 在 $\xi$ 处的严格局部化}.
\end{enonce*}

由 \hyperref[VIII.prop3.9]{3.9}, 纤维 $\Oo_{{X,}_{\xi}}$ 可表示为
\begin{equation*}
\Oo_{{X,}_{\xi}}=\varinjlim\Gamma(X',\Oo_{X'}){\quoi},
\tag{$\ast$}
\end{equation*}
其中 $X'$ 遍历由 $\xi$ 标点的平展 $X$-概形的反向滤过范畴,
即配有一个 $X$-态射 $\xi\to X'$ 的那些概形.
由归纳极限的传递性, 可以把右端的 $\Gamma(X',\Oo_{X'})$
换成 $\Oo_{X',x'}$, 其中 $x'$ 是 $\xi$ 在 $X'$ 中的像,
从而得到
\begin{equation*}
\Oo_{{X,}_{\xi}}=\varinjlim_i A_i
\tag{\ast\ast}
\end{equation*}
其中 $A_i$ 遍历 $A=\Oo_{X,x}$ 上本质有限型且平展的代数,
并配有一个 $k$-同态 $k(A_i)\to\Omega$.
(这里 $x$ 表示 $\xi$ 在 $X$ 中的像.)
$A_i$ 之间的过渡同态当然是使相应三角图
\[
\xymatrix@R=1.25cm@C=1.75cm{
k(A_i)  \ar[r] \ar[rd]& k(A_j)\ar[d]    \\
                      &  \Omega.
}
\]
交换的局部 $A$-代数同态 $A_i\to A_j$.
这意味着, 若从 $A_i$ 到 $A_j$ 存在一个允许的同态, 它就是唯一的.
因此归纳极限 $({\ast}{\ast})$ 是在一个向上滤过的有序指标集上的极限.

\setcounter{subsection}{3}
\subsection{}\etiquette[4.4]{VIII.subsec4.4}

表示式\pageoriginale $({\ast}{\ast})$ 表明,
$X$ 在 $\xi$ 处的严格局部环实质上仅依赖于通常的局部环
$A=\Oo_{X,x}$ ($x=u(\xi)$) 和 $k=k(A)$ 的扩张 $\Omega$.
它也称为 $A$ 的\emph{严格 Hensel 化}
(相对于 $k$ 的给定可分闭扩张 $\Omega$), 记为 $A^{hs}$.
关于 $A^{hs}$ 的一般性质, 参见 \EGA IV 18.8.
我们只指出, 这个构造再次表明 $A^{hs}$ 在 $A$ 上平坦,
且若 $A$ 是 Noether 环, 它也是.
此外容易证明 (\loccit), 局部同态
$$
A\to A^{hs}
$$
连同 $k$-同态
$$
k(A^{hs})\to\Omega
$$
解决了如下泛问题:
对于给定的 $A$ 和 $k$ 的扩张 $\Omega$, 寻找局部同态
$$
A\to A'{\quoi},
$$
其中 $A'$ \emph{严格局部}, 并配有一个 $k$-同态
$$
k(A')\to\Omega{\quoi}.
$$
特别地, $A^{hs}$ \emph{确实是一个严格局部环}
(这就说明了 \hyperref[VIII.def4.3]{4.3} 中术语的合理性).

\subsection{}\etiquette[4.5]{VIII.subsec4.5}

沿用 \hyperref[VIII.def4.3]{4.3} 的记号,
并选取 $x$ 的一个仿射开邻域 $U$.
注意, 在 \hyperref[VIII.prop3.9]{3.9} 中,
显然可以只考虑在 $U$ 上仿射的 $X'$.
它们构成一个仿射概形的伪滤过逆系,
所以我们\pageoriginale 处于 \hyperrefext[VII][VII.sec5]{5}
的一般情形 (见 VI 5.12 b).
由 $\Oo_{X,\xi}$ 的表示式 $(\ast)$, 有一个 $X$-同构
\begin{equation*}
\varprojlim_{\substack{X'\ \mbox{\rm 平展}\\ \mbox{在 }U\mbox{ 上仿射},\\
X'_{\xi}\mbox{\rm -标点}}}X'\simeq\Spec\Oo_{X,\xi}
\tag{\ast\ast\ast}
\end{equation*}
它把 $\Spec(\Oo_{X,\xi})$ 作为“\emph{$\xi$ 的平展邻域的极限}”
这一直观的几何意义准确化了.

\begin{enonce*}{命题 4.6}\etiquette[4.6]{VIII.prop4.6}
设 $X$ 是一个严格局部概形, $x$ 是它的闭点,
因此也是 $X$ 的一个几何点.
则函子 $F\mapsto\Gamma(X,F):\widetilde{X_{\text{\rm ét}}}\to(\Ens)$
典范同构于纤维函子 $F\mapsto F_x$,
因而与任意 $\varinjlim$ 交换\footnote{不要求滤过!}.
\end{enonce*}

事实上, 根据定义 \hyperref[VIII.def4.2]{4.2} (ii),
\hyperref[VIII.prop3.9]{3.9} 中考虑的由 $x$ 标点的平展 $X$-概形范畴的终对象 $X$
支配每个其他对象;
所以 $\varinjlim F(X')$ 典范同构于 $F(X)$, 正如所需.

特别地, $\Gamma$ 在 Abel 层范畴上诱导的函子正合, 因此:

\begin{enonce*}{推论 4.7}\etiquette[4.7]{VIII.cor4.7}
在 \hyperref[VIII.prop4.6]{4.6} 的假设下,
对 $X$ 上的每个 Abel 层 $F$, 有
$$
H^q(X,F)=0\quad\mbox{\rm 当}\quad q\neq0{\quoi}.
$$
\end{enonce*}

\begin{enonce*}{推论 4.8}\etiquette[4.8]{VIII.cor4.8}
设 $X$ 是一个概形, $\xi$ 是 $X$ 的一个几何点,
$\bar X=\Spec\Oo_{X,\xi}$ 是相应的严格局部化,
$F$ 是 $X$ 上的一个层, $\bar F$ 是它在 $\bar X$ 上的逆像.
则有函子同构
$$
F_\xi\simeq\Gamma(\bar X,\bar F){\quoi}.
$$
\end{enonce*}

事实上,\pageoriginale $\xi$ 也可视为 $\bar X$ 的一个几何点,
只须应用 \hyperref[VIII.prop4.6]{4.6} 和纤维的传递性
(\hyperref[VIII.subsec3.4]{3.4}) 即可.


% END INCLUDED SGA4_English_current_buildable_source_20260906/source/08/drafts/SGA4_VIII_LINES_577_804_English.texfrag


% BEGIN INCLUDED SGA4_English_current_buildable_source_20260906/source/08/drafts/SGA4_VIII_LINES_805_1011_English.texfrag
\section{在计算 $R^qf_*$ 的纤维中的应用}\etiquette[5]{VIII.sec5}

\subsection{}\etiquette[5.1]{VIII.subsec5.1}

设
$$
f:X\to Y
$$
是一个概形态射, $F$ 是 $X$ 上的一个 Abel 层.
我们希望确定 $R^qf_*(F)$.
根据 \hyperref[VIII.thm3.5]{3.5},
这几乎就等于知道对 $y\in Y$ 的几何纤维 $R^qf_*(F)_{\bar y}$.
现在 (\hyperrefext[V][V.prop5.1]{5.1}),
$R^qf_*(F)$ 是 $Y$ 上的预层
$$
\Hh^q:Y'\longmapsto H^q(\fprod{X}{Y'}{Y},F)
$$
的相伴层, 因而 (\hyperref[VIII.prop3.9]{3.9})
$$
R^qf_*(F)_{\bar y}=\varinjlim_{Y'}\Hh^q(Y'){\quoi},
$$
其中归纳极限取遍由 $\bar{y}$ 标点的平展 $Y$-概形的反向滤过范畴.
选定 $y$ 的一个仿射开邻域 $U$ 后,
可以把右端的极限限制在那些仿射且位于 $U$ 上的 $Y'$ 所成的共尾范畴上.
于是得到
\begin{equation*}
R^qf_*(F)_{\bar y}\simeq
\varinjlim H^q(\fprod{X}{Y'}{Y},F){\quoi},
\tag{5.1.1}
\end{equation*}\etiquette[5.1.1]{VIII.eq5.1.1}
其中 $Y'$ 在上述范畴中变化.

置
$$
\bar Y=\Spec(\Oo_{Y,\bar y})=\varprojlim Y'
$$
(见 \hyperref[VIII.subsec4.5]{4.5}), 以及
$$
\bar X=\fprod{X}{\bar Y}{Y}{\quoi}.
$$
注意,\pageoriginale 在逆系 $\fprod{X}{Y'}{Y}$ 中
(它由 $Y'$ 所成逆系经基变换 $X\to Y$ 得到), 过渡态射都是仿射的.
因此我们处于 \hyperrefext[VII][VII.subsec5.1]{5.1} 的一般情形,
\loccit 中 $\varprojlim$ 的构造立即给出一个典范同构
$$
\bar X=\varprojlim\fprod{X}{Y'}{Y}{\quoi}.
$$
用 $\bar F$ 表示 $F$ 在 $\bar X$ 上的逆像层.
于是得到一个典范同态
$$
\varinjlim_{Y'}H^q(\fprod{X}{Y'}{Y},F)
\longrightarrow H^q(\bar X,\bar F),
$$
因而与 \eqhyperref[VIII.eq5.1.1]{5.1.1} 比较, 得到一个典范同态
\begin{equation*}
R^qf_*(F)_{\bar y}\longrightarrow H^q(\bar X,\bar F)
\tag{5.1.2}{\quoi},
\end{equation*}\etiquette[5.1.2]{VIII.eq5.1.2}
它显然关于 $F$ 具有函子性.

现在假设 $f:X\to Y$ 拟紧且拟分离.
则 $\fprod{X}{Y'}{Y}\to Y'$ 也如此;
由于 $Y'$ 仿射, 概形 $\fprod{X}{Y'}{Y}$ 拟紧且拟分离.
应用 \eqhyperref[VIII.eq5.1.1]{5.1.1} 及极限过渡定理
\hyperrefext[VII][VII.cor5.8]{5.8}, 得到:

\begin{enonce*}{定理 5.2}\etiquette[5.2]{VIII.thm5.2}
设 $f:X\to Y$ 是一个拟紧且拟分离的概形态射,
$F$ 是 $X$ 上的一个 Abel 层, $y$ 是 $Y$ 的一个点,
$\bar{y}$ 是 $y$ 上方与 $k(y)$ 的某个可分闭包 $k(\bar y)$ 对应的几何点.
设 $\bar Y=\Spec(\Oo_{Y,\bar y})$ 是相应的严格局部化,
$\bar X=\fprod{X}{\bar Y}{Y}$,
$\bar F$ 是 $F$ 在 $\bar X$ 上的逆像.
则典范同态 \eqhyperref[VIII.eq5.1.2]{5.1.2} 是一个同构:
$$
R^qf_*(F)_{\bar y}\simeq H^q(\bar X,\bar F){\quoi}.
$$
\end{enonce*}

这个陈述实际上把确定层\pageoriginale $R^qf_*(F)$ 的纤维的问题,
归结为确定严格局部概形上的某个概形的上同调的问题,
下文将不断使用它.
从技术上说, \hyperref[VIII.thm5.2]{5.2}
解释了 Hensel 环和严格局部环在研究平展上同调中的重要作用.

\begin{enonce*}[remark]{注 5.3}\etiquette[5.3]{VIII.rem5.3}
对于任意集合层, 这个陈述对 $R^0f_*(F)=f_*(F)$ 仍成立.
此外, 当 $F$ 是群层 (不一定交换) 时,
采用 $R^1f_*(F)$ 的通常定义 (见 Giraud 的学位论文),
\hyperref[VIII.thm5.2]{5.2} 对 $R^1f_*(F)$ 也仍成立.
\end{enonce*}

\setcounter{subsection}{3}
\subsection{}\etiquette[5.4]{VIII.subsec5.4}

现在假设 $f$ 是一个\emph{有限}态射.
则 $\bar f:\bar X\to\bar Y$ 有限,
由于 $\bar Y$ 严格局部, 可知 $\bar X$ 是严格局部概形 $\bar X_i$ 的有限不交并.
因此, 利用 \hyperref[VIII.cor4.7]{4.7}, 得到
\begin{equation*}
H^q(\bar X,\bar F)=0\quad\mbox{当}\quad q\neq0{\quoi}.
\tag{5.4.1}
\end{equation*}\etiquette[5.4.1]{VIII.eq5.4.1}
此外, $\bar X$ 的分支 $\bar X_i$ 对应于 $\bar X$ 中位于 $\bar Y$
的闭点 $\bar{y}$ 上方的点 $\bar x_i$,
即 $\bar X_{\bar y}=X_y\otimes_{k(y)}k(\bar y)$ 的点.
这些点也可视为 $X$ 的几何点,
此时 $\bar X_i$ 恰好是 $X$ 在 $\bar x_i$ 处的严格局部化.
因此 (\hyperref[VIII.cor4.8]{4.8})
$$
H^0(\bar X_i,\bar F)\simeq F_{\bar x_i}{\quoi},
$$
从而
\begin{equation*}
H^0(\bar X,\bar F)=\mathop{\prod}_iF_{\bar x_i}{\quoi}.
\tag{5.4.2}
\end{equation*}\etiquette[5.4.2]{VIII.eq5.4.2}
这是\pageoriginale 一个关于集合层 $F$ 具有函子性的同构
(这里不必限制于 Abel 层).
结合 \hyperref[VIII.thm5.2]{5.2} 和 \hyperref[VIII.rem5.3]{5.3},
公式 \eqhyperref[VIII.eq5.4.1]{5.4.1} 和 \eqhyperref[VIII.eq5.4.2]{5.4.2} 给出:

\begin{enonce*}{命题 5.5}\etiquette[5.5]{VIII.prop5.5}
设 $f:X\to Y$ 是一个有限概形态射, $y$ 是 $Y$ 的一个点.
则对 $X$ 上的每个层 $F$, 有一个关于 $F$ 具有函子性的同构
$$
f_*(F)_{\bar y}\simeq
\prod_{\bar x\in X_y\otimes_{k(y)}k(\bar y)}F_{\bar x}{\quoi}.
$$
因此, $f_*(F)$ 的形成与每个基变换 $Y'\to Y$ 交换;
若 $F$ 是 Abel 层, 则
$$
R^qf_*(F)=0\quad\mbox{若}\quad q>0{\quoi}.
$$
\end{enonce*}

注意, \hyperref[VIII.prop5.5]{5.5} 的第一个公式
实际上独立于 \hyperref[VIII.thm5.2]{5.2}
和一般的极限过渡定理 \hyperrefext[VII][VII.thm5.7]{5.7};
它与 \hyperref[VIII.thm3.5]{3.5} 一起表明,
$F\mapsto f_*(F)$ 在 Abel 层范畴上是一个\emph{正合函子},
从而再次得到, 当 $q\neq0$ 时 $R^qf_*(F)=0$.
下面是 \hyperref[VIII.prop5.5]{5.5} 的一个小变式:

\begin{enonce*}{推论 5.6}\etiquette[5.6]{VIII.cor5.6}
设 $f:X\to Y$ 是一个整态射.
则对 $X$ 上的每个 Abel 层 $F$, 有
$$
R^qf_*(F)=0\quad\mbox{当}\quad q\neq0{\quoi}.
$$
此外, 集合层上的函子 $f_*$ 与每个基变换 $Y'\to Y$ 交换.
\end{enonce*}

事实上, 由 \hyperref[VIII.thm5.2]{5.2},
只须证明当 $Y$ 严格局部时, 对 $X$ 上的每个 Abel 层 $F$, 都有 $H^q(X,F)=0$.
此时 $Y=\Spec A$, $X=\Spec B$, 其中 $B$ 是一个整 $A$-代数;
写成\pageoriginale $B=\varinjlim B_i$,
其中 $B_i$ 遍历 $B$ 的有限型子代数 (它们事实上在 $A$ 上有限),
便有 $X=\varprojlim_iX_i$, 其中 $X_i=\Spec B_i$.
根据 \hyperrefext[VII][VII.cor5.13]{5.13},
只须证明当 $q\neq0$ 时 $R^qf_*(F_i)=0$,
其中 $F_{i}$ 是 $X_{i}$ 上的一个 Abel 层;
这由 \hyperref[VIII.prop5.5]{5.5} 得出.
\hyperref[VIII.cor5.6]{5.6} 的最后一个断言,
以同样方式归结到 \hyperref[VIII.prop5.5]{5.5} 而得证.

\begin{enonce*}{推论 5.7}\etiquette[5.7]{VIII.cor5.7}
设 $f:X\to Y$ 是一个浸入. 则对 $X$ 上的每个层 $F$, 典范态射
$$
f^*f_*(F)\to F
$$
是一个同构.
\end{enonce*}

由于 $f$ 可分解为开浸入和闭浸入的复合,
而 \hyperref[VIII.cor5.7]{5.7} 对开浸入显然成立 (IV 第 3 节),
只须考虑 $f$ 是闭浸入的情形.
须证明对每个 $x\in X$, 相应同态
$$
(f^*f_*(F))_{\bar x}\to F_{\bar x}
$$
是双射. 由纤维的传递性 (\hyperref[VIII.subsec3.4]{3.4}),
左端恰是纤维 $f_*(F)_{\bar x}$;
因此只须验证典范同态
$$
f_*(F)_{\bar x}\to F_{\bar x}
$$
是双射, 这直接由 \hyperref[VIII.prop5.5]{5.5} 得出.

\begin{enonce*}[remark]{注 5.8}\etiquette[5.8]{VIII.rem5.8}
利用 \hyperref[VIII.prop5.5]{5.5},
并再次像 \hyperref[VIII.cor5.6]{5.6} 中那样论证, 容易证明:
若 $f:X\to Y$ 是一个整态射, $F$ 是 $X$ 上的一个群层 (不一定交换),
则 $R^1f_*(F)$ 是 $Y$ 上的终层 (比较注 \hyperref[VIII.rem5.3]{5.3}).
\end{enonce*}


% END INCLUDED SGA4_English_current_buildable_source_20260906/source/08/drafts/SGA4_VIII_LINES_805_1011_English.texfrag


% BEGIN INCLUDED SGA4_English_current_buildable_source_20260906/source/08/drafts/SGA4_VIII_LINES_1012_1556_English.texfrag
\section{支集}\pageoriginale\etiquette[6]{VIII.sec6}

设 $U$ 是概形 $X$ 的一个 Zariski 开子集.
则 $U \in \Ob X_{\text{\rm ét}}$,
而且 $U$ 实际上是 $X_{\text{\rm ét}}$ 的终对象 $X$ 的一个子对象,
从而定义了 $\widetilde{X_{\text{\rm ét}}}$ 的终对象的一个子对象 $\widetilde{U}$,
\ie $X$ 的平展拓扑斯 $X_{\text{\rm ét}}$ 的一个“开子集”
(IV \hyperref[IV.def8.3]{8.3}).

\begin{enonce*}{命题 6.1}\etiquette[6.1]{VIII.prop6.1}
上述映射 $U \mapsto \overset{\sim}{U}$
是从 $X$ 的 (Zariski 意义下的) 开子集所成的有序集,
到平展拓扑斯 $\widetilde{X_{\text{\rm ét}}}$ 的开子集所成集合的一个同构.
\end{enonce*}

由于 $X_{\text{\rm ét}} \to \widetilde{X_{\text{\rm ét}}}$ 全忠实,
立即可知映射 $U \mapsto \widetilde{U}$ 保持序结构,
\ie $(U \subset V) \Leftrightarrow (\widetilde{U} \subset \widetilde{V})$;
特别地, 上述映射是单射.
为证明它是满射, 考虑终层 $\widetilde{X}$ 的一个子层 $F$,
并考虑 $X_{\text{ét}/ F}$ 的对象,
\ie $X$ 上满足 $F(X') \neq ∅$ 的平展概形 $X'$.
由于 $X' \to X$ 平展, 它是开映射 (\EGA IV 2.4.6);
特别地, 它的 (集合意义下的) 像 $\text{Im}(X')$ 是开的.
设 $U$ 是所有 $\text{Im}(X')$ $(X'\in \Ob X_{\text{ét} /  F})$ 的并所成的开子集.
由于族 $X'\to U$ $(X'\in \Ob X_{\text{ét}/F})$ 是满射的,
因而是覆盖族, 可得 $U \in \Ob X_{\text{ét}/F}$,
所以 $ X_{\text{ét}/F} =  X_{\text{ét}/U}$,
从而 $F = U$, \cqfd.

根据 \hyperref[VIII.prop6.1]{6.1},
此后可以无歧义地谈论 $X$ 的“开子集”,
无须说明这是通常 Zariski 意义下的概念, 还是平展拓扑意义下的概念.

\begin{enonce*}{推论 6.2}\etiquette[6.2]{VIII.cor6.2}
设 $U$ 是 $X$ 的一个开子集, $j : U \to X$ 是典范浸入.
则函子
$$
j^{\ast} :  \widetilde{X_{\text{\rm ét}}} \to \widetilde{U_{\text{\rm ét}}}
$$
诱导\pageoriginale 一个范畴等价
$$
\widetilde{X_{\text{\rm ét}}} / \widetilde{U} \to
\widetilde{U_{\text{\rm ét}}}{\quoi}.
$$
\end{enonce*}

事实上, 直接可验证:
对每个存在有限 $\varprojlim$ 的景 $C$,
及 $C$ 的终对象 $e$ 的每个子对象 $U$,
若考虑由 $j(S) = S \times U$ 定义的函子 $j : C \to C_{/U}$,
则 $j$ 是一个景态射,
且 $j^{\ast} : \widetilde{C} \to \widetilde{C/U}$
诱导一个等价 $\widetilde{C}/\widetilde{U} \to \widetilde{C / U}$.
将这个一般事实与 $U_{\text{ét}}$ 典范同构于 $X_{\text{ét}/U}$ 这一事实结合起来即可.

形象地说, 可以将 \hyperref[VIII.cor6.2]{6.2} 表述为:
一方面在概形的通常意义下、另一方面在拓扑斯的意义下,
“限制到开子集”的运算彼此相容.
下面是“限制到闭子集”运算的类似相容性:

\begin{enonce*}{定理 6.3}\etiquette[6.3]{VIII.thm6.3}
设 $X$ 是一个概形, $Y$ 是 $X$ 的一个闭子概形,
$U = X - Y$ 带有诱导结构,
$i : Y \to X$ 和 $j : U \to X$ 是典范浸入.
则函子
$$
i_{\ast} : \widetilde{Y_{\text{\rm ét}}} \to \widetilde{X_{\text{\rm ét}}}
$$
全忠实, 且若 $F \in \Ob \widetilde{X_{\text{\rm ét}}}$,
则 $F$ 同构于形如 $i_{\ast} (G)$ 的层,
当且仅当 $j^{\ast}(F)$ 同构于 $U$ 上的终层 $\widetilde{U}$.
\end{enonce*}

证明. 由于 $i_{\ast}$ 和 $i^{\ast}$ 互相伴随,
$i_{\ast}$ 全忠实也等价于函子同态
$$
i^{\ast} (i_{\ast}(G)) \to G
$$
是一个\pageoriginale 同构, 这正是 \hyperref[VIII.cor5.7]{5.7}.
另一方面, 若 $G \in \Ob \widetilde{Y_{\text{ét}}}$,
则利用 \hyperref[VIII.cor6.2]{6.2}, 显然可验证
$j^{\ast} (i_{\ast}(G))$ 是 $U$ 上的终层.
反过来, 若 $F \in \Ob \widetilde{X_{\text{ét}}}$
满足 $j^{\ast}(F)$ 是终层,
我们证明 $F$ 同构于形如 $i_{\ast}(G)$ 的层;
等价地, 证明典范同态
$$
F \to i_{\ast} i^{\ast} F
$$
是一个同构.
仍然只须对每个 $x \in X$,
验证在 $\bar{x}$ 处的几何纤维上诱导的同态是同构.
当 $x \in U$ 时, 这恰好是对 $F$ 所作的假设.
当 $x \in Y$ 时, 由纤维的传递性,
归结为验证由
$$
i^{\ast} (F) \to i^{\ast} (i_{\ast} i^{\ast} (F))
$$
在 $\bar{x}$ 处的纤维上诱导的同态是同构;
但将 \hyperref[VIII.cor5.7]{5.7} 应用于 $F = i^{\ast} F$ 和 $i$,
可知上述同态是同构.
这就完成了 \hyperref[VIII.thm6.3]{6.3} 的证明.

\begin{enonce*}{推论 6.4}\etiquette[6.4]{VIII.cor6.4}
函子 $i_{\ast}$ 诱导 $Y$ 上的 Abel 层范畴,
到 $X$ 上限制在 $U = X - Y$ 上为零的 Abel 层所成范畴的一个等价.
\end{enonce*}

因此, 按照通常的用法,
我们常将 $Y$ 上的一个 Abel 层与 $X$ 上的一个在 $U$ 上为零的 Abel 层等同.

\setcounter{subsection}{4}
\subsection{}\etiquette[6.5]{VIII.subsec6.5}

根据 \hyperref[VIII.prop6.1]{6.1},
当 $X$ 是概形时, 我们知道如何将拓扑斯
$ℰ = \widetilde{X_{\text{ét}}}$ 的“开子集”解释为
$X$ 的通常意义下的开子集 $U$;
而根据 \hyperref[VIII.thm6.3]{6.3},
在这一等同之下, IV 3 中的“剩余拓扑斯”\pageoriginale $ℰ_\complement U$
典范等价于拓扑斯 $\widetilde{Y_{\text{ét}}}$,
其中 $Y = X - U$ 带有任意一种使 $Y$ 成为概形的诱导结构.
因此可以将 IV 3 的结果应用于此,
特别是 IV \hyperref[IV.theo9.5.4]{3.3} 中
用三元组 $(F', F'', u)$ 描述 $X$ 上的层 $F$ 的结果,
其中 $F'$ 是 $Y$ 上的层, $F''$ 是 $U$ 上的层,
$u$ 是从 $F'$ 到 $i^{\ast} j_{\ast} (F'')$ 的同态.
下文将自由使用 IV.3 中的记号 $i^{!}, j_{!}$, 它们表示函子
\begin{align*}
j_{!} : \widetilde{(U_{\text{ét}})}_{\ab} &\to
\widetilde{(X_{\text{ét}})}_{\ab}{\quoi},\\
i^{!} : \widetilde{(X_{\text{ét}})}_{\ab} &\to
\widetilde{(Y_{\text{ét}})}_{\ab}{\quoi},
\end{align*}
其中下标“$\ab$”表示 Abel 层范畴.
这些函子给出 IV \hyperref[IV.prop14.6]{3.7} 中的两个正合列.

\subsection{}\etiquette[6.6]{VIII.subsec6.6}

根据前面的讨论以及 IV \hyperref[IV.subsec8.5]{8.5} 中引入的一般术语,
对 $X$ 上的 Abel 层 $F$ 或这种层的一个截面 $\varphi$,
适宜引入 \emph{$F$ 的支集}, \resp \emph{$\varphi$ 的支集} 这一概念:
它是所考虑对象为零的最大开子集的补集,
也就是一个 (通常意义下, \ie Zariski 意义下的) 闭子集.

在当前情形中, 还应将一般记号
\hyperrefext[V][V.prop4.3]{4.3} $H^q (ℰ_{\complement U} F)$,
$ℋ^q_{\complement U} (F)$ 换成记号 $H^{q}_Y (X, F)$,
$ℋ^{q}_Y (F)$;
前者表示一个 Abel 群, 后者表示 $Y$ 上的一个 Abel 层
(或等价地, $X$ 上的一个在 $U$ 上为零的 Abel 层).
因此有 $\mathcal H^q_Y(F)=R^q i^!(F)$, \etc

关于上述函子的一般性质, 参见 \hyperrefext[V][V.sec6]{6}.

\section{纤维函子的特殊化态射}\pageoriginale
\etiquette[7]{VIII.sec7}

\subsection{}\etiquette[7.1]{VIII.subsec7.1}

在第 ${\rm N^{\circ}}$ 4 节中,
我们已看到如何对概形 $X$ 的每个几何点 $\xi$ 配上一个严格局部 $X$-概形
$$
\bar{X}(\xi) = \Spec \Oo_{X, \xi}{\quoi},
$$
它实际上只依赖于 $\xi$ 的像 $x$ 上方的那个几何点,
后者由 $k(x)$ 在 $\Omega = k(\xi)$ 中的可分闭包 $\bar{k(x)}$ 定义.
下文常常仅考虑在 $X$ 上\emph{代数可分}的几何点 $\xi = \Spec \Omega$,
\ie $\Omega$ 是 $k(x)$ 的一个可分闭包, \ie $\xi = \bar{x}$.
若一个 $X$-概形 $Z$ 与形如 $\Spec \Oo_{X, \bar{x}}$ 的概形 $X$-同构,
就称它是 \emph{$X$ 的一个严格局部化}; 于是可见
$$
\xi \mapsto \Spec \sheaf{\Oo}_{X, \xi} = \bar{X}(\xi)
$$
是从 $X$ 上\emph{代数可分}的几何点范畴到 $X$ 的严格局部化所成的
$X$-概形范畴的一个等价,
这里两个范畴中均仅取同构作为态射.

若取所有 $X$-态射作为态射, 最后一个断言就不再正确,
因为在第二个范畴中可能有不是同构的 $X$-态射.
于是作如下定义:

\begin{enonce*}{定义 7.2}\etiquette[7.2]{VIII.def7.2}
设 $\xi$, $\xi'$ 是概形 $X$ 的两个几何点.
从 $\xi'$ 到 $\xi$ 的特殊化箭头,
是相应的严格局部化之间的任意一个 $X$-态射
$$
\bar{X} (\xi') \to \bar{X} (\xi).
$$
若存在从 $\xi'$ 到 $\xi$ 的特殊化箭头,
就\pageoriginale 称 $\xi$ 是 $\xi'$ 的一个特殊化,
或称 $\xi'$ 是 $\xi$ 的一个一般化.
\end{enonce*}

注意, 特殊化箭头按显然的方式复合,
所以以特殊化箭头为态射, \emph{$X$ 的几何点构成一个范畴},
它等价于 $(\Sch)_{/X}$ 中由 $X$ 的严格局部化组成的全子范畴
(由 $X$ 上代数可分的几何点组成的全子范畴也与之等价).

\begin{enonce*}{引理 7.3}\etiquette[7.3]{VIII.lem7.3}
设 $X$ 是一个概形, $Z$ 是一个 $X$-概形,
同构于平展 $X$-概形 $X_{i}$ 的一个伪滤过投射极限,
其过渡态射仿射 \emph{(\hyperrefext[VII][VII.subsec5.1]{5.1})};
设 $\xi'$ 是 $X$ 的一个几何点,
$Z' = \bar{X} (\xi')$ 是相应的严格局部化. 则:
\begin{enumerate}
\renewcommand{\theenumi}{\alph{enumi}}
\renewcommand{\labelenumi}{\theenumi)}
\item 限制映射
$$
\Hom_X (Z', Z) \to \Hom_X(\xi', Z)
$$
是双射.

\item 两端非空的充要条件是 $\xi'$ 在 $X$ 中的像 $x'$ 属于 $Z$ 的像.

\item 设 $T$ 是一个 $Z$-概形.
$T$ 是 $Z$ 的一个严格局部化的充要条件是它为 $X$ 的一个严格局部化.
\end{enumerate}
\end{enonce*}

证明.
\begin{enumerate}
\renewcommand{\theenumi}{\alph{enumi}}
\renewcommand{\labelenumi}{\theenumi)}
\item 直接归结到 $Z$ 是某个 $X_{i}$ 的情形,
\ie $Z$ 在 $X$ 上平展.
再经基变换 $Z'\to X$, 可假设 $Z'\xrightarrow{\sim} X$,
于是结论由 \hyperref[VIII.def4.2]{4.2} (ii) 得出.

\item 注意, 若 $x'$ 是 $Z$ 的某个点 $z$ 的像,
则 $k(z)$ 必是 $k(x')$ 的一个可分代数扩张;
因此存在从后者到 $k(\xi')$ 的一个 $k(x')$-同态, 从而得出结论.

\item 证明\pageoriginale 是一道容易的练习, 留给读者.
\end{enumerate}

从 \hyperref[VIII.def7.2]{7.2} 和 \hyperref[VIII.lem7.3]{7.3} 得到:

\begin{enonce*}{命题 7.4}\etiquette[7.4]{VIII.prop7.4}
设 $\xi$, $\xi'$ 是概形 $X$ 的两个几何点.
则限制映射给出从 $\xi'$ 到 $\xi$ 的特殊化箭头所成的集合
$\Hom_{X} (\bar{X}(\xi'), \bar{X}(\xi))$,
到从 $\xi'$ 到 $\bar{X}(\xi)$ 的 $X$-态射所成集合的一个双射.
\end{enonce*}

\begin{enonce*}{推论 7.5}\etiquette[7.5]{VIII.cor7.5}
$\xi$ 是 $\xi'$ 的一个特殊化的充要条件,
是 $\xi$, $\xi'$ 在 $X$ 中的像 $x$, $x'$ 也如此,
\ie $x$ 属于 $\{x'\}$ 的闭包.
\end{enonce*}

由于 $\Spec \Oo_{X, \xi} \to \Spec \Oo_{X, x}$ 忠实平坦
(\hyperref[VIII.subsec4.4]{4.4}), 因而是满射,
所以只须应用 \hyperref[VIII.lem7.3]{7.3} 的第二个断言.

\begin{enonce*}{推论 7.6}\etiquette[7.6]{VIII.cor7.6}
对每个概形 $X$,
设 $\sheaf{\Pt}(X)$ 为 $X$ 上的几何点范畴
(或等价地, 为 $X$ 上代数可分的几何点范畴), 其态射是特殊化箭头.
设 $\xi$ 是概形 $X$ 的一个几何点,
$\bar{X}(\xi)$ 是相应的严格局部化.
则有一个范畴等价
$$
\sheaf{\Pt}(\bar{X}(\xi)) \xrightarrow{\sim} \sheaf{\Pt}(X)_{/\xi}{\quoi},
$$
它将 $\bar{X}(\xi)$ 的每个几何点 $\xi'$
对应到相应的 $X$ 上的几何点,
再配以由结构态射 $\xi' \to \bar{X}(\xi)$
按 \hyperref[VIII.prop7.4]{7.4} 得出的到 $\xi$ 的特殊化态射.
\end{enonce*}

换言之, \emph{给定几何点 $\xi$ 的一个一般化 $\xi'$,
实质上等价于给定 $\bar{X}(\xi)$ 的一个几何点}.
此外注意, 由 \hyperref[VIII.lem7.3]{7.3} c),
相应同态\pageoriginale $\bar{X}(\xi') \to \bar{X}(\xi)$
使 $\bar{X}(\xi')$ 成为 $\bar{X}(\xi)$ 关于 $\xi'$ 的严格局部化.

\setcounter{subsection}{6}
\subsection{}\etiquette[7.7]{VIII.subsec7.7}

现在, 对 $X$ 的每个几何点 $\xi$ 和 $X$ 上的每个层 $F$,
将纤维 $F_{\xi}$ 解释为 $\Gamma (\bar{X} (\xi), \bar{F}(\xi))$,
其中 $\bar{F}(\xi)$ 是 $F$ 在 $\bar{X}(\xi)$ 上的逆像
(\hyperref[VIII.cor4.8]{4.8}).
于是可见, 每个特殊化箭头
$$
u : \xi'\to \xi
$$
都诱导一个关于 $F$ 具有函子性的同态:
$$
u^{\ast} : F_{\xi} \to F_{\xi'}{\quoi},
$$
称为与特殊化箭头 $u$ 相应的\emph{特殊化同态}.
由层的逆像的传递性显然可知,
对于一个复合特殊化箭头, 有:
$$
(wu)^{\ast} = u^{\ast} w^{\ast}{\quoi}.
$$

\subsection{}\etiquette[7.8]{VIII.subsec7.8}

% SGA-SRC-000768: the scan supplies no exact citation; reconstructed from the mathematics.
若 $ℰ$ 是一个拓扑斯, 回顾第 IV 讲 6.1 和 6.2:
拓扑斯 $ℰ$ 的一个“纤维函子”, 或 (滥用语言) 一个“点”,
是从“终拓扑斯” $(\Ens)$
(它同构于退化为一个点的空间上的层范畴!) 到 $ℰ$ 的任意一个态射,
\ie 任意一个与任意归纳\sisi{极限}{极限}以及有限投射\sisi{极限}{极限}交换的函子
$$
\varphi : ℰ \to (\Ens)
$$
应将 $ℰ$ 的纤维函子所成的集合视为
$\Hhom(ℰ, (\Ens))$ 的一个全子范畴的对象集,
这个子范畴称为拓扑斯 $ℰ$ 的\emph{纤维函子范畴},
其反范畴称为 \emph{$ℰ$ 的点范畴}, 记为 $\sheaf{\Pt}(ℰ)$,
\cf (IV \hyperref[IV.def6.1]{6.1}).

当 $ℰ$ 形如 $\widetilde{X_{\text{ét}}}$, 其中 $X$ 是概形时,
我们已在\pageoriginale \hyperref[VIII.def3.3]{3.3}
和 \hyperref[VIII.subsec7.7]{7.7} 中定义了一个函子
\begin{equation*}
\sheaf{\Pt} (X) \to \sheaf{\Pt} (\widetilde{X_{\text{ét}}}){\quoi},\tag{$\ast$}
\end{equation*}
其中左端在 \hyperref[VIII.cor7.6]{7.6} 中定义.
据此, 有如下定理:

\begin{enonce*}{定理 7.9}\etiquette[7.9]{VIII.thm7.9}
设 $X$ 是一个概形.
上述函子 $(\ast)$ 是从 $X$ 上的几何点范畴
(以特殊化箭头为态射),
到平展拓扑斯 $\widetilde{X_{\text{\rm ét}}}$ 的点范畴
($\widetilde{X_{\text{\rm ét}}}$ 上的纤维函子范畴的\emph{反范畴}) 的一个等价.
\end{enonce*}

由于本讨论班后面不会用到这个定理,
这里仅给出证明概要, 并且对宇宙问题允许作某些宽松处理.
\begin{enumerate}
\renewcommand{\theenumi}{\alph{enumi}}
\renewcommand{\labelenumi}{\theenumi)}
\item \emph{所考虑的函子全忠实}.
对每个 $X \in \Ob X_{\text{ét}}$,
设 $\Xv : \widetilde{X_{\text{ét}}} \to (\Ens)$ 由
$$
\Xv (F) = \Hom(X, F) = F(X){\quoi}.
$$
定义. 注意, 纤维函子
$$
\mathcal{E}_{\xi} : F \mapsto F_{\xi}
$$
可写为
$$
\mathcal{E}_{\xi} \simeq \varinjlim_i X^{\vee}_i{\quoi},
$$
其中 $X_{i}$ 是 $X$ 上的仿射平展概形,
以某个滤过范畴为指标
(\cf \hyperref[VIII.def4.3]{4.3} 和 \hyperref[VIII.subsec4.5]{4.5}).
因此对每个函子 $\varphi : \widetilde{X_{\text{ét}}} \to (\Ens)$, 有:
$$
\Hom(\mathcal{E}_{\xi}, \varphi) \simeq \varprojlim_{i}
\Hom(X^{\vee}_i, \varphi){\quoi},
$$
另一方面, 众所周知有一个双函子同构
$$
\Hom(X^{\vee}_i, \varphi) \simeq \varphi (X_i){\quoi}.
$$

当\pageoriginale $\varphi$ 形如 $\Ee_{\xi'}$, 从而
$$
\varphi \simeq \varinjlim_j {X'}_j^{\vee}{\quoi},
$$
时, 便有一个自然双射
\begin{equation*}
\Hom (\mathcal{E}_{\xi}, \mathcal{E}_{\xi'}) \simeq
\varprojlim_{i} \varinjlim_{j} \Hom_X (X'_j, X_i){\quoi}.\tag{$\ast$}
\end{equation*}
另一方面, 在概形范畴中有
$$
\bar{X}(\xi) = \varprojlim_{i} X_i,\quad \bar{X} (\xi') =
\varprojlim_{j} X'_j{\quoi},
$$
从而
$$
\Hom_{X} (\bar{X} (\xi'), \bar{X}(\xi)) \simeq \varprojlim_{i}
\Hom_X (\bar{X} (\xi'), X_i){\quoi},
$$
又因为 $X_i$ 在 $X$ 上局部有限表示, 有
$$
\Hom_X (\bar{X} (\xi'), X_i) \simeq \varinjlim_{j}
\Hom_X (X'_j, X_i){\quoi},
$$
从而
\begin{equation*}
\Hom_X (\bar{X} (\xi'), \bar{X} (\xi) ) \simeq
\varprojlim_{i} \varinjlim_{j} \Hom_X (X'_j,
X_i){\quoi}.\tag{${\ast}{\ast}$}
\end{equation*}
比较 $(\ast)$ 和 $({\ast}{\ast})$ 即得所需结论
(还须验证相容性, 这留给读者).

\item \emph{所考虑的函子本质满射}.
\end{enumerate}

\setcounter{subsection}{9}
\subsubsection{}\etiquette[7.9.1]{VIII.subsubsec7.9.1}

首先注意, 若 $C$ 是一个有限 $\varprojlim$ 可表的景,
$\widetilde{C} = ℰ$ 是相应拓扑斯,
则 $ℰ$ 上的每个纤维函子 $\varphi$ 都可表示为形如
$$
\check{Y} (F) = F(Y) = \Hom(\widetilde{Y}, F), Y \in \Ob C{\quoi},
$$
的函子的一个“滤过归纳极限”
(其中 $\widetilde{Y}$ 是与 $Y$ 相伴的层), 方法如下
\footnote{这是 IV \hyperref[IV.eq6.8.3]{6.8.3} 的一个特例.}.
考虑由二元组 $(Y, \xi)$ 组成的范畴\pageoriginale $C_{/\varphi}$,
其中 $Y \in \Ob C$, $\xi \in \varphi (\widetilde{Y})$
(态射按显然的方式定义), 并注意有一个显然的同态
\begin{equation*}
\varinjlim_{(Y, \xi)\in C^0_{/\varphi}} \check{Y} \to
\varphi\tag{$\ast$}
\end{equation*}
因为
$$
\Hom (\varinjlim_{C^0_{/\varphi}} \check{Y},\varphi) \simeq
\varprojlim_{C_{/ \varphi}} \Hom(\check{Y}, \varphi) \simeq
\varprojlim_{C_{/ \varphi}} \varphi (\widetilde{Y}){\quoi}.
$$
$\varprojlim_{C_{/ \varphi}} \varphi(\widetilde{Y}),$
中有一个典范元素,
它由将每个 $(Y, \xi)$ 对应到 $\varphi (\widetilde{Y})$ 中的元素 $\xi$ 而得.
利用 $C$ 中存在有限 $\varprojlim$,
以及函子 $Y \mapsto \varphi(\widetilde{Y})$ 与它们交换这两个事实,
可见 $C_{/\varphi}$ 中存在有限 $\varprojlim$;
因此 $C^0_{/\varphi}$ 更是滤过的.
利用每个层 $F$ 都是形如 $\widetilde{Y}$ 的层的归纳极限,
以及 $\varphi$ 与归纳极限交换这两个事实,
容易得到函子同态 $(\ast)$ 是双射,
从而给出所述表示.

\subsubsection{}\etiquette[7.9.2]{VIII.subsubsec7.9.2}

还要注意, 若 $ℰ$ 是一个拓扑斯, $Y$ 是 $ℰ$ 的一个对象,
从而定义一个拓扑斯 $ℰ_{/Y}$,
则给定 $ℰ_{/Y}$ 的一个纤维函子 $\varphi$,
等价于给定一个二元组 $(\varphi, \xi)$,
其中 $\varphi$ 是 $ℰ$ 的一个纤维函子,
$\xi \in \varphi(Y)$.
与纤维函子 $\psi : ℰ_{/Y} \to (\Ens)$ 相应的二元组为 $(\varphi, \xi)$,
其中 $\varphi$ 是复合 $Z \mapsto \varphi(Z) = \psi (Z \times Y)$,
$\xi \in \varphi(Y) = \psi(Y \times Y)$
则是 $\psi(Y)$ 的唯一元素在 $Y \times Y$ 的对角态射下的像.

\subsubsection{}\etiquette[7.9.3]{VIII.subsubsec7.9.3}

现在回到景 $C = X_{\text{ét}}$ 的情形,
设 $\varphi : \widetilde{X_{\text{ét}}} \to (\Ens)$
是 $X$ 的平展拓扑斯上的一个纤维函子.
考虑\pageoriginale $\varphi$ 在 $\widetilde{X_{\text{ét}}}$ 的如下全子范畴上的限制:
由这个拓扑斯的开子集组成的子范畴,
或等价地 (\hyperref[VIII.prop6.1]{6.1}),
由 $X$ 的开子集 $U$ 组成的子范畴.
由此得到一个与任意上确界及有限下确界交换的映射
$$
\varphi_0 : \Uu (X) \to \{0, 1\}
$$
(注意, 对 $U \in \Uu (X)$, $\varphi(U)$ 或者为空,
或者仅由一点组成, 按前一种或后一种情形, 分别令 $\varphi_{0} (U) = 0$ 或 1.)

利用 $X$ 的每个不可约闭子集恰有一个泛点这一事实
\footnote{\ie 用 IV \hyperref[IV.subsubsec4.2.1]{4.2.1} 的术语来说,
$X$ 是\emph{清醒的}.
所述断言是 IV \hyperref[IV.subsubsec4.2.3]{4.2.3} 的一个特例.},
容易看出, $\varphi_0$ 由唯一的 $x \in X$ 通过条件
$$
\varphi_0 (U) = 1\quad \Longleftrightarrow\quad x \in U {\quoi},
$$
定义, \ie
$$
\varphi (U) \ne\varnothing \Leftrightarrow  x \in U {\quoi}.
$$
设 $F\in \Ob \widetilde{X_{\text{ét}}}$.
则 $F$ 在终层 $X$ 中的像是一个开子集;
由于
$$
\varphi (F) \to \varphi(U)
$$
是满射, 可见 $\varphi(F) \neq ∅$,
当且仅当 $\varphi(U) \neq ∅$, \ie 当且仅当 $x \in U$.
设 $C_{/\varphi}$ 是由二元组 $(X', \xi)$ 组成的范畴,
其中 $X'\in \Ob X_{\text{ét}}$, $\xi \in \varphi (X')$.
\hyperref[VIII.subsubsec7.9.1]{7.9.1} 中已经指出,
$C^{0}_{/ \varphi}$ 滤过且连通.
此外, 由 \hyperref[VIII.subsubsec7.9.2]{7.9.2},
并将上述讨论应用于 $X'$ 而非 $X$,
若 $(X', \xi) \in \Ob C_{/\varphi}$,
则与它相应的有唯一一点 $x' \in X'$,
满足对平展态射 $X'' \to X'$,
$\xi$ 属于 $\varphi (X'') \to \varphi(X')$ 的像,
当且仅当 $x'$ 属于 $X''$ 的像.
由此可知, 在由 $(X', \xi)$ 组成的范畴 $C$ 上,
这些 $X'$ 的 $\varprojlim$ 概形存在, 并且是 $X$ 的一个严格局部化 $Z$,
从而对应于一个几何点 $\xi$ 和一个纤维函子 $\mathcal{E}_{\xi}$.
于是对\pageoriginale 每个层 $F$, 有
$$
\mathcal{E}_{\xi} (F) \simeq \varinjlim_{C^0_{/\varphi}} F(X'){\quoi}.
$$
根据 \hyperref[VIII.subsubsec7.9.1]{7.9.1},
可得 $\varphi$ 同构于 $\mathcal{E}_{\xi}$,
这就完成了 \hyperref[VIII.thm7.9]{7.9} 的证明.

% END INCLUDED SGA4_English_current_buildable_source_20260906/source/08/drafts/SGA4_VIII_LINES_1012_1556_English.texfrag


% BEGIN INCLUDED SGA4_English_current_buildable_source_20260906/source/08/drafts/SGA4_VIII_LINES_1557_1976_English.texfrag
\section{整态射的两个谱序列}\etiquette[8]{VIII.sec8}

\begin{enonce*}{命题 8.1}\etiquette[8.1]{VIII.prop8.1}
设 $f : X \to Y$ 是一个满射整态射, $F$ 是 $Y$ 上的一个 Abel 层.
对每个整数 $i$, 设 $\Hh^{i} (F)$ 是 $(\Sch)_{/Y}$ 上由
$$
\Hh^{i} (F) (Z) = H^i (Z, F_Z){\quoi},
$$
定义的预层, 其中 $F_Z$ 是 $F$ 在 $Z$ 上的逆像.
则存在一个 (关于 $F$ 具有函子性的) 谱序列
$$
H^{\ast} (X, F) \Longleftarrow E^{pq}_2 = H^p (X /Y, \Hh^q (F))
$$
(“下降谱序列”).
\end{enonce*}

当然, 对 $(\Sch)_{/Y}$ 上的预层 $G$,
符号 $H^p (X / Y, G)$ 表示第 $p$ 个相对 \v{C}ech 上同调群,
它用 $C^n(X / Y, G) = G(X^{n+1})$ 所成的复形定义,
其中 $X^m$ 表示 $(\Sch)_{/Y}$ 中的 $m$ 次幂.
为证明 \hyperref[VIII.prop8.1]{8.1},
设 $p_n : X^{n+1} \to Y$ 为投射, 并置
$$
A^n = p_{n\ast} (𝒵_{X^{n + 1}}){\quoi},
$$
其中 $𝒵_{X^{n + 1}}$ 表示 $X^{n+1}$ 上的常值层 $𝒵$.
则 $A^n$ 是 $Y$ 上的一个单纯 Abel 层的各个分量,
从而组成 $Y$ 上的一个 Abel 层复形 $A^{\ast}$.
注意有一个显然的同态
\begin{equation*}
𝒵_Y \to A^{0}{\quoi}.\tag{$\ast$}
\end{equation*}
据此\pageoriginale, 有如下引理:

\begin{enonce*}{引理 8.2}\etiquette[8.2]{VIII.lem8.2}
配备同态 $(\ast)$ 的复形 $(A^{\ast})$ 是 $𝒵_Y$ 的一个消解.
更一般地, 对 $Y$ 上的每个 Abel 层 $F$,
$A^{\ast} \otimes_{𝒵}F$ 是 $F$ 的一个消解.
\end{enonce*}

\noindent
证明: 可以假设 $Y$ 仿射,
从而 $Y = \Spec A$, $X = \Spec B$, 其中 $B$ 是一个整 $A$-代数.
有 $B = \varinjlim_i B_i$, 其中 $B_i$ 遍历 $B$ 的有限生成子代数,
它们因而是有限的;
从而 $X=\varprojlim X_{i}$, 其中 $X_{i}=\Spec B_{i}$.
利用 \hyperrefext[VII][VII.cor5.11]{5.11},
可见增广复形 $A^{\ast}(X/Y)$ 是增广复形 $A^{\ast}(X_{i}/Y)$ 的归纳极限.
这将我们归结到 $f$ \emph{有限}的情形.
只须证明对 $Y$ 的每个几何点 $\bar{y}$,
复形 $A^{\ast}_{\bar y}$ 是 $𝒵_{\bar{y}}$ 的消解,
并且与某个 $F$ 作张量积后也如此.
但若 $X_{\bar{y}}$ 是 $X$ 在 $\bar{y}$ 处的纤维,
由 \hyperref[VIII.prop5.5]{5.5} 可知,
复形 $A^{\ast}_{\bar{y}}$ 正是类似的复形 $A^{\ast}(X_{\bar{y}}/\bar{y})$.
这就归结到 $Y$ 是可分闭域 $k$ 的谱的情形.
利用 \hyperref[VIII.exem1.3]{1.3} a) 和 b),
甚至可假设 $k$ 代数闭且 $X$ 约化,
从而 $X$ 形如 $I_{Y}$, 其中 $I$ 是有限集.
此时复形 $A^{\ast}\otimes F$ 可视为指标集 $I$ 上以 $F_{\bar{y}}$ 为系数的平凡上链复形,
因此确实是 $F_{\bar{y}}$ 的一个消解.

于是, 作为 \hyperref[VIII.lem8.2]{8.2} 的推论, 得到一个谱序列
$$
H^{\ast}(Y,F)\Longleftarrow E^{pq}_{2}=H^{p}(H^{q}(Y,A^{\ast}\otimes
F)){\quoi},
$$
只须将初始项写明. 有一个典范同态\pageoriginale
$$
A^{n}\otimes F\longrightarrow p_{n_{\ast}}p^{\ast}_{n}(F){\quoi},
$$
它是同构, 这一点仍可归结到 $f$ 有限的情形, 再过渡到纤维来验证.
利用 \hyperref[VIII.prop5.5]{5.5}, 得到
$$
H^{q}(Y,A^{n}\otimes F)\simeq
H^{q}(X^{n+1},p^{\ast}_{n}(F))=
\Hh^{q}(F)(X^{n+1}){\quoi},
$$
这确实给出了 \hyperref[VIII.prop8.1]{8.1} 中所述的初始项.

\begin{enonce*}[remark]{注 8.3}\etiquette[8.3]{VIII.rem8.3}
\begin{enumerate}
\renewcommand{\theenumi}{\alph{enumi}}
\renewcommand{\labelenumi}{\theenumi)}
\item 若给定 $Y$ 的一个由闭子集 $Y_{i}$ 组成的\emph{局部有限}覆盖,
并令 $X=\coprod_{i}Y_{i}$,
则典范态射 $f:X\rightarrow Y$ 有限,
而谱序列 \hyperref[VIII.prop8.1]{8.1} 的初始项可写为
由 $H^{q}(Y_{i_{0}\ldots i_{p}},F|Y_{i_{0}\ldots i_{p}})$
所定义的复形的上同调,
其中 $Y_{i_{0}\ldots i_{p}}=Y_{i_{0}}\cap\ldots\cap Y_{i_{p}}$.
因此它是通常拓扑空间的局部有限闭覆盖的 Leray 谱序列
(TF 第 II 章 5.2.4) 的类似物.
而且后者也可按 \hyperref[VIII.prop8.1]{8.1} 的方法,
推广为一个“有限”态射 (\ie 有限纤维的固有态射) 的谱序列.

\item 注意 \hyperref[VIII.prop8.1]{8.1} 的谱序列,
与 $Y$ 的覆盖 $(X_{i})$ 的 Leray 谱序列之间的类似性
(在后一种情形中, $X_i$ 在 $Y$ 上平展);
形式地, 后者可对 $X=\coprod_{i}X_{i}$
写出谱序列 \hyperref[VIII.prop8.1]{8.1} 而得到.
事实上, 这两个谱序列很可能具有一个共同推广,
只要给定的态射族 $(X_{i}\rightarrow Y)$
对于变动的基概形\pageoriginale 上的平展层所成的纤维化范畴,
是一个“普遍有效下降族”, 这个推广就有效
(\cf 下文第 ${\rm n^{\circ}}$ 9 节).
在通常拓扑中也有类似问题:
对分别假设为开覆盖、或局部有限闭覆盖的两类覆盖 Leray 谱序列,
寻找一个共同推广
\footnote[1]{在这些文字写成之后,
P.~Deligne 发展了“下降谱序列”的一般理论,
\cf \Exp V bis.}.
\end{enumerate}
\end{enonce*}

\begin{enonce*}{命题 8.4}\etiquette[8.4]{VIII.prop8.4}
设 $Y$ 是一个概形, $\pi$ 是一个投射有限群,
$(\pi_{i})_{i}$ 是 $\pi$ 的有限离散商群组成的投射系,
$(X_{i})_{i\in I}$ 是 $Y$ 的一个主覆盖投射系,
其群为 $\pi_{i}$, 过渡同态 $X_{j}\rightarrow X_{i}$
与算子群上的同态 $\pi_{j}\rightarrow \pi_{i}$ 相容;
设 $X=\varprojlim X_{i}$
\emph{(\cf \hyperrefext[VII][VII.subsec5.1]{5.1})},
从而群 $\pi$ 作用在 $Y$-概形 $X$ 上.
设 $F$ 是 $Y$ 上的一个 Abel 层.
则存在一个“\emph{Hochschild--Serre}”谱序列
(关于 $F$ 具有函子性)
$$
H^{\ast}(Y,F)\Longleftarrow
E^{pq}_{2}=H^{p}(\pi,\varinjlim_{i}H^{q}(X_{i},F_{X_{i}})){\quoi},
$$
其中 $F_{X_{i}}$ 是 $F$ 在 $X_i$ 上的逆像,
$H^{p}(\pi,-)$ 表示 Galois 上同调.
\end{enonce*}

这里写出的项 $E^{pq}_{2}$ 还
(由 $H^{p}(\pi,-)$ 的定义及归纳极限的传递性) 同构于
$$
E^{pq}_{2}=\varinjlim_{i}
H^{p}(\pi_{i},H^{q}(X_{i},F_{X_{i}})){\quoi},
$$
这表明, 只须找到一个 (依赖于指标 $i$ 的) 谱序列归纳系
$$
H^{\ast}(Y,F)\Longleftarrow
{}^{i}E^{pq}_{2}=H^{p}(\pi_{i},H^{q}(X_{i},F_{X_{i}})){\quoi}.
$$
这将我们归结为在 $\pi$ 有限的情形中定义谱序列.
此时 $f:X\rightarrow Y$ 是 $Y_{\text{\rm ét}}$ 中的一个覆盖态射,
因此可以\pageoriginale 写出这个态射的 Leray 谱序列.
根据典范同构
$$
X^{n}\simeq X\mathop×^{\pi}\pi^{n}{\quoi},
$$
一个众所周知的计算表明:
对 $X_{\text{\rm ét}}$ 上每个把和变为积的预层 $\Hh$,
$C^{\ast}(X/Y,\Hh)$ 正是 $\pi$ 的以 $\Hh(X)$ 为系数的齐次上链复形,
由此得出所述初始项的形式.

这个证明有一个避免任何计算的变式:
% SGA-SRC-000769: corrupt source notation reconstructed as a contracted product.
将 $E\mapsto X\mathop{\times}^{\pi}E$
视为从有限 $\pi$-集合景到 $Y$ 的平展景的一个态射,
并写出这个态射的 Leray 谱序列.
最后, 当 $\pi$ 有限时,
也可以将 Hochschild--Serre 谱序列视为
\hyperref[VIII.prop8.1]{8.1} 对态射 $f:X\rightarrow Y$ 的一个特例.

\begin{enonce*}{推论 8.5}\etiquette[8.5]{VIII.cor8.5}
假设 $Y$ 拟紧且拟分离.
则 \hyperref[VIII.prop8.4]{8.4} 的谱序列具有形式
$$
H^{\ast}(Y,F)\Longleftarrow E^{pq}_{2}=H^{p}(\pi,
H^{q}(X,F_{X})){\quoi}.
$$
\end{enonce*}

事实上, 由 \hyperrefext[VII][VII.cor5.8]{5.8}, 此时有典范同构
$$
\varinjlim_{i} H^{q}(X_{i},F_{X_{i}})\simeq H^{q}(X,F_{X}){\quoi}.
$$

\begin{enonce*}{推论 8.6}\etiquette[8.6]{VIII.cor8.6}
设 $Y$ 是一个局部 Hensel 概形, 闭点为 $y$,
$F$ 是 $Y$ 上的一个 Abel 层,
$F_{0}$ 是 $Y_{0}=\Spec (k(y))$ 上的诱导层.
则典范同态
$$
H^{n}(Y,F)\longrightarrow H^{n}(Y_{0},F_{0})
$$
都是\pageoriginale 同构.
\end{enonce*}

事实上, 设 $\bar y$ 是 $y$ 上方的一个几何点,
对应于 $k(y)$ 的一个可分闭包 $\bar{k(y)}=k(\bar{y})$;
由于 $Y$ 是 Hensel 的,
$Y$ 在 $\bar y$ 处的严格局部化 $X$
是由 $\bar y$ 标点的连通平展 Galois 覆盖的投射极限,
将这些覆盖记为 $X_i$,
于是 \hyperref[VIII.cor8.5]{8.5} 的假设适用.
由于根据 \hyperref[VIII.cor4.7]{4.7},
当 $q>0$ 时 $H^{q}(X,F_{X})=0$, 得到同构
$$
H^{n}(Y,F)\xleftarrow{\sim} H^{n}(\pi, F(X)){\quoi},
$$
其中 $\pi$ 是 $X$ 在 $Y$ 上的“Galois 群”,
同构于 $\bar{k(y)}$ 在 $k=k(y)$ 上的 Galois 群.
类似地 (或根据 \hyperref[VIII.cor2.3]{2.3}), 得到
$$
H^{n}(Y_{0},F_{0})\xleftarrow{\sim} H^{n}(\pi,F_{0}(X_{0})){\quoi},
$$
其中 $X_{0}=\fprod{X}{Y}{Y_{0}}\simeq \Spec (k(\bar{y}))$.
但由 \hyperref[VIII.cor4.8]{4.8},
限制同态 $F(X)\rightarrow F_{0}(X_{0})$ 是同构,
由此立即得出 \hyperref[VIII.cor8.6]{8.6} 的结论.

\section{平展层的下降}\etiquette[9]{VIII.sec9}

本节在本讨论班后面不会用到, 初读时可以略去.

\begin{enonce*}{命题 9.1}\etiquette[9.1]{VIII.prop9.1}
设\pageoriginale $f:X'\rightarrow X$ 是一个满射
(\resp 普遍淹没 \emph{(\SGA I IX 2.1)}) 概形态射.
则函子 $f^{\ast}:X_{\text{\rm ét}}\rightarrow X'_{\text{\rm ét}}$
忠实且“保守” \emph{(\cf \hyperref[VIII.cor3.6]{3.6})}
(\resp 诱导一个从 $X_{\text{\rm ét}}$ 上的层范畴,
到配备相对于 $f:X'\rightarrow X$ 的下降数据的
$X'_{\text{\rm ét}}$ 上的层范畴的\emph{全忠实}函子).
\end{enonce*}

第一个断言直接由纤维函子的传递性
(\hyperref[VIII.subsec3.4]{3.4}) 以及 (\hyperref[VIII.cor3.7]{3.7}) 得出;
第二个断言 (也可表述为: 对于变动概形上的平展层所成的纤维化范畴,
$f$ 是一个\emph{下降态射})
也就是说, 对 $X$ 上的两个层 $F$, $G$, 自然图
\begin{equation*}
\Hom(F,G)\rightarrow \Hom (F',G')\rightrightarrows \Hom (F'',G'')\tag{$\ast$}
\end{equation*}
正合, 其中 $F'$ 和 $G'$ (\resp $F''$ 和 $G''$)
是 $F$ 和 $G$ 在 $X'$ 上
(\resp 在 $X''=\fprod{X'}{X'}{X}$ 上) 的逆像.
取 $F$ 为终层, 这个断言给出以下推论:

\begin{enonce*}{推论 9.2}\etiquette[9.2]{VIII.cor9.2}
设 $f:X'\rightarrow X$ 是一个普遍淹没态射.
则对 $X$ 上的每个层 $F$, 图
$$
\Gamma(X,F)\rightarrow \Gamma(X',F')\rightrightarrows \Gamma(X'',F'')
$$
正合, 其中 $X''=\fprod{X'}{X'}{X}$,
$F'$, $F''$ 是 $F$ 在 $X'$, $X''$ 上的逆像.
\end{enonce*}

我们用以下引理证明 \hyperref[VIII.prop9.1]{9.1}:

\begin{enonce*}{引理 9.3}\etiquette[9.3]{VIII.lem9.3}
假设 $f$ 普遍淹没.
设 $F$ 是 $X$ 上的一个层,
用 $S(F)$ 表示 \sisi{$X$}{$F$} 的子层所成的集合,\pageoriginale
并类似地定义 $S(F')$, $S(F'')$.
则由自然映射组成的图
$$
S(F)\rightarrow S(F')\rightrightarrows S(F'')
$$
正合.
\end{enonce*}

\noindent
证明. $S(F)\rightarrow S(F')$ 是单射,
直接由函子 $f^{\ast}$ 保守这一事实得出;
事实上, 若 $F_i$ $(i=1,2)$ 是 $F$ 的两个子层,
满足 $f^{\ast}(F_{1})=f^{\ast}(F_{2})$,
则包含态射 $F_{3}=F_{1}\cap F_{2}\rightarrow F_{i}$ $(i=1,2)$ 都是同构,
因为应用函子 $f^{\ast}$ 后它们成为同构, 从而 $F_{1}=F_{2}$.
只须证明: 若 $G'$ 是 $F'$ 的一个子层,
满足 $\pr^{\ast}_{1}(G')=\pr^{\ast}_{2}(G')$,
则 $G'$ 是 $F$ 的某个子层的逆像.
注意, 可在 $X_{\text{\rm ét}}$ 中找到一个满态射 $F_{1}\rightarrow F$,
其中 $F_{1}$ 可由 $X$ 上的一个平展概形\emph{代表}.
引入 $F_{2}=\fprod{F_{1}}{F_{1}}{F}$,
并类似地引入 $F'_{1}$, $F''_{1}$, $F'_{2}$, $F''_{2}$,
这些层给出一个集合图
\[
\xymatrix@R=.75cm{
  S(F)  \ar[r] \ar[d]& S(F') \ar[r]\ar[d] & S(F'') \ar[d]\\
S(F_1) \ar[r] \ar@<3pt>[d] \ar@<-3pt>[d]& S(F'_1) \ar@<3pt>[r]
  \ar@<-3pt>[r] \ar@<3pt>[d] \ar@<-3pt>[d]& S(F''_{1})\ar@<3pt>[d] \ar@<-3pt>[d] \\
S(F_2)  \ar[r] & S(F'_2)\ar@<3pt>[r]
  \ar@<-3pt>[r] & S(F''_2),
}
\]
其中各列正合, 这分别来自 $X$, $X'$, $X''$ 上的层范畴的正合性性质.
一次标准的图追踪表明,
要证明第一行正合, 只须证明第 2 行和第 3 行正合.
对于第 2 行, 这来自 $F_{1}$ 可表、
$X'\rightarrow X$ 普遍\pageoriginale 淹没以及 \SGA\sisi{}{ 1} IX 2.3.
另一方面, $F_{2}$ 也可表,
因为它是 $\fprod{F_{1}}{F_{1}}{X}$ 的一个子层,
后者可表, 因而可应用 \hyperref[VIII.prop6.1]{6.1}.
所以第 3 行也正合, \cqfd

现在证明 \hyperref[VIII.prop9.1]{9.1},
\ie 对每个与 $F'$, $G'$ 上的下降数据相容的态射 $u':F'\rightarrow G'$,
该态射 $u'$ 都来自一个态射 $u:F\rightarrow G$.
设 $H=F\times G$, 从而 $H'=F'\times G'$, $H''=F''\times G''$;
则 $u'$ 的\emph{图}是 $H'$ 的一个子层 $\Gamma'$,
它的两个逆像等于同一个态射 $u'':F''\rightarrow G''$ 的图.
因此由 \hyperref[VIII.lem9.3]{9.3},
$\Gamma'$ 来自 $H$ 的一个子层 $\Gamma$.
我们断言, $\Gamma$ 是某个态射 $u:F\rightarrow G$ 的图,
\ie 由 $\pr_{1}$ 诱导的态射 $p:\Gamma\rightarrow F$ 是同构:
事实上, 它经过基变换 $X'\rightarrow X$ 后成为同构,
然后应用 \hyperref[VIII.prop9.1]{9.1} 中已经证明的部分即可.
于是态射 $u:F\rightarrow G$ 具有所需性质, \cqfd

\begin{enonce*}{定理 9.4}\etiquette[9.4]{VIII.thm9.4}
设 $f:X'\rightarrow X$ 是一个满射概形态射.
假设以下条件之一成立:
\begin{enumerate}
\renewcommand{\theenumi}{\alph{enumi}}
\renewcommand{\labelenumi}{\theenumi)}
\item $f$ 整.

\item $f$ 固有.

\item $f$ 平坦且局部有限表示
\footnote[1]{这个第二条件很可能实际上是多余的.}.

\item $X$ 离散 (例如某个域的谱).
\end{enumerate}

则对于 $(\Sch)$ 上由变动概形上的平展层组成的纤维化范畴 $\Ff$,
$f$ 是一个\emph{有效下降态射},
\ie 函子 $f^{\ast}$ 诱导 $X$ 上的层范畴\pageoriginale
与配备相对于 $f:X'\rightarrow X$ 的下降数据的 $X'$ 上的层范畴之间的一个\emph{等价}.
\end{enonce*}

\setcounter{subsection}{4}
\setcounter{subsubsection}{0}
\subsubsection{}\etiquette[9.4.1]{VIII.subsubsec9.4.1}

设 $F'$ 是 $X'$ 上的一个层,
配有相对于 $f:X'\rightarrow X$ 的下降数据.
用公式
$$
G(Y)=\Ker (F'(Y')\rightrightarrows F''(Y'')){\quoi}.
$$
定义一个函子
$$
G:X^{0}_{\text{\rm ét}}\longrightarrow (\Ens)
$$
立即可见 $G$ 是平展拓扑的一个\emph{层}.
另一方面, 有一个典范单同态 $G\rightarrow f_{\ast}(F')$,
从而有一个同态 $f^{\ast}(G)\rightarrow F'$,
后者显然与下降数据相容.
只须判断这个同态是否为\emph{同构}.
注意, 也可用下图的正合性来定义 $G$:
$$
G\rightarrow f_{\ast}(F')\rightrightarrows g_{\ast}(F''){\quoi},
$$
所以, 对每个\emph{与 $f_{\ast}(F')$ 和 $g_{\ast}(F'')$ 的形成交换的}
基变换 $Y\xrightarrow{h}X$,
$h^{\ast}(G)=G_{Y}$ 可视为由 $F'_{Y'}$
经 $Y'\rightarrow Y$ “下降”所得的层;
当然, 同态 $f_{Y}{}^{\ast}(G_{Y})\rightarrow (F_{Y})'$
正是从 $f^{\ast}(G)\rightarrow F'$ 经基变换所得的同态.
现在\emph{假设基变换} $f:Y=X'\rightarrow X$
\emph{与 $f_{\ast}(F')$ 和 $g_{\ast}(F'')$ 的形成交换},
并注意, 由于 $Y'=\fprod{X'}{X'}{X}$ 在 $Y=X'$ 上有一个截面,
$Y'\rightarrow Y$ 对于每个纤维化范畴都是\emph{有效}下降态射,
尤其对 $\mathcal{F}$ 如此;
因此 $f^{\ast}_{Y}(G_{Y})\rightarrow (F_{Y})'$ 是同构.
于是 $f^{\ast}(G)\rightarrow F'$ 在基变换
$Y'=\fprod{X'}{X'}{X}\rightarrow X'$ 后成为同构.
由于后者有一个截面, $f^{\ast}(G)\rightarrow F'$ 是同构,
所以所考虑的 $F'$ 上的下降数据\pageoriginale 有效.

\subsubsection{}\etiquette[9.4.2]{VIII.subsubsec9.4.2}

由 \hyperref[VIII.cor5.6]{5.6},
这就在情形 a) 中证明了定理 \hyperref[VIII.thm9.4]{9.4}.
情形 b) 也来自以下事实:
若 $f$ 固有, 则 $f_{\ast}$ 与\emph{每个}基变换 $Y'\rightarrow Y$ 交换
(这一点将在后面的一讲中证明 (XII \hyperref[XII.thm5.1]{5.1} (i))).

\subsubsection{}\etiquette[9.4.3]{VIII.subsubsec9.4.3}

在情形 c) 中, 由于问题在 $X$ 上是局部的,
可以假设 $X$ 仿射, 并且存在一个在 $X$ 上有限表示、拟有限且忠实平坦的仿射概形 $X'_{1}$,
再配以一个 $X$-态射 $X'_{1}\rightarrow X'$,
这将我们归结到 $f$ 拟有限的情形.
再在 $X$ 上局部化, 这次是在平展拓扑意义下,
可见存在 $X'$ 的一个开子集 $X'_{1}$,
使 $X'_{1}\rightarrow X$ 有限且满射,
从而归结到 $f$ 有限且满射的情形, 它已在 a) 中处理过.

\subsubsection{}\etiquette[9.4.4]{VIII.subsubsec9.4.4}
在情形 d) 中, 可以 (在 $X$ 上局部化) 假设 $X$ 仅由一个点组成,
根据 \hyperref[VIII.thm1.1]{1.1}, 甚至可以假设它是某个域 $k$ 的谱.
此外, $f$ 普遍开 (\EGA IV 2.4.9),
所以 \hyperref[VIII.prop9.1]{9.1} 的假设成立,
而 \hyperref[VIII.subsubsec9.4.1]{9.4.1} 的论证再次适用,
只须经 $Y=\Spec k$ 作基变换, 其中对某个 $x\in X$, 有 $k=k(x)$.
但此时 $f_{\ast}$ 仍与基变换 $Y\rightarrow X$ 交换,
正如我们以后将看到的那样
(XVI \hyperref[XVI.compl1.4]{1.4} 和 \hyperref[XVI.cor1.5]{1.5}).

这就完成了 \hyperref[VIII.thm9.4]{9.4} 的证明.


\ifx\danslelivre\undefined
\end{document}
\fi

% END INCLUDED SGA4_English_current_buildable_source_20260906/source/08/drafts/SGA4_VIII_LINES_1557_1976_English.texfrag


